% Minimal preamble for Open Logic Project sources (pinned commit 1e960beff9ed7835bf3e3f1335e21af3439cd107).
%
% OLP's own master preamble cannot run on the project host: TeX Live 2025 there lacks
% prooftree.sty, ntheorem.sty, extarrows.sty, mathdots.sty, isomath.sty and subfiles.sty.
% Following the project's own rule ("if the author preamble needs packages the toolchain lacks,
% do not invent macros -- use a minimal working preamble"), this file declares, verbatim in
% spirit and numbering behaviour, exactly the theorem-like environments OLP defines in
% sty/open-logic-book-envs.sty (thm/lem/prop/cor/defn/ex/prob/rem/note/case/axiom/conv),
% plus the shorthand accents OLP uses (`!!a`, `!!b`, `!!{...}`).
%
% Author text is untouched; only the packaging is ours.
\documentclass[10pt,a4paper]{article}

\usepackage{amsmath}
\usepackage{amssymb}
\usepackage{amsthm}
\usepackage{graphicx}
\usepackage{hyperref}
\usepackage{enumitem}
\usepackage[framemethod=default]{mdframed}

% ---- OLP shorthand accents (from sty/open-logic.sty) ----
% NB: \! is taken by LaTeX (negative thin space), so we only define the !!a / !!b families
% and the \olup helper used by them.
\newcommand*{\olup}[1]{\MakeUppercase{#1}}
\def\!!!a#1{\olup{#1}}
\def\!a#1{\olup{#1}}
\def\!b#1{\olup{#1}}

% ---- theorem-like environments, same family numbering as OLP ----
\theoremstyle{plain}
\newtheorem{thm}{Theorem}[section]
\newtheorem{lem}[thm]{Lemma}
\newtheorem{prop}[thm]{Proposition}
\newtheorem{cor}[thm]{Corollary}
\newtheorem{axiom}{Axiom}

\theoremstyle{definition}
\newtheorem{defn}[thm]{Definition}
\newtheorem{ex}[thm]{Example}
\newtheorem{prob}{Problem}[section]
\newtheorem{probd}{Problem}
\newtheorem{conv}[thm]{Convention}

\theoremstyle{remark}
\newtheorem{rem}[thm]{Remark}
\newtheorem{note}[thm]{Note}
\newtheorem{case}[thm]{Case}

% ---- environments OLP defines as wrappers ----
\newenvironment{probtag}[1]{\begin{prob}}{\end{prob}}
\newenvironment{explain}{}{}
\newenvironment{explain*}{}{}
\newenvironment{editorial}{}{}
\newenvironment{digress}{}{}
\newenvironment{digress*}{}{}
\newenvironment{history}{}{}
\newenvironment{reference}{}{}
\newcommand*{\olfileid}[3]{}
\newcommand*{\olsection}[1]{\section*{#1}}
\newcommand*{\olchapter}[3]{}
\newcommand*{\olpart}[2]{}
\newcommand*{\olimport}[1]{}
\newcommand*{\olinclude}[1]{}
\newcommand*{\olpath}[1]{}
\newcommand*{\ollangid}[1]{}
\newcommand*{\problemtag}[1]{}
\newcommand*{\problemauthor}[1]{}
\newcommand*{\ollabel}[1]{\label{#1}}
% \olref takes 0-3 OPTIONAL arguments plus one mandatory label:
%     \olref{lab}  |  \olref[from]{lab}  |  \olref[from][to]{lab}  |  \olref[a][b][c]{lab}
% A single-argument definition silently swallows "[sth]" and emits a broken reference
% (measured on the host: the text came out as "??iso]x"). The chain below consumes the
% brackets and cross-references only the mandatory argument.
\makeatletter
\newcommand{\olref}{\@ifnextchar[{\olref@i}{\olref@final}}
\def\olref@i[#1]{\@ifnextchar[{\olref@ii}{\olref@final}}
\def\olref@ii[#1]{\@ifnextchar[{\olref@iii}{\olref@final}}
\def\olref@iii[#1]{\olref@final}
\def\olref@final#1{\ref{#1}}
\makeatother
\newcommand*{\olnewpage}{\clearpage}

% ===================== OLP compatibility layer (1) =====================

\makeatletter

% ---------------- proof trees (OLP's prooftree is a bussproofs derivative) ----------------
\providecommand{\AxiomC}[1]{}
\providecommand{\UnaryInfC}[1]{}
\providecommand{\BinaryInfC}[1]{}
\providecommand{\TrinaryInfC}[1]{}
\providecommand{\QuaternaryInfC}[1]{}
\providecommand{\QuinaryInfC}[1]{}
\providecommand{\RightLabel}[1]{}
\providecommand{\LeftLabel}[1]{}
\providecommand{\doubleLine}{}
\providecommand{\noLine}{}
\providecommand{\DisplayProof}{}
\providecommand{\fCenter}{}
\newenvironment{prooftree}{}{}
\newenvironment{ltprooftree}{}{}
\newenvironment{oltableau}{}{}

% ---------------- logical notation ----------------
\providecommand{\lif}{\mathbin{\rightarrow}}
\providecommand{\lforall}{\mathord{\forall}}
\providecommand{\lexists}{\mathord{\exists}}
\providecommand{\ltrue}{\top}
\providecommand{\lfalse}{\bot}
\providecommand{\Ltrue}{\top}
\providecommand{\Lfalse}{\bot}
\providecommand{\Sequent}{\mathrel{\Rightarrow}}
\providecommand{\Proves}{\mathrel{\vdash}}
\providecommand{\Entails}{\mathrel{\models}}
\providecommand{\Subst}[3]{#1[#2/#3]}
\providecommand{\sFmla}[2]{#1(#2)}
\providecommand{\pFmla}[2]{#1(#2)}
\providecommand{\formula}{\relax}
\providecommand{\True}{\top}
\providecommand{\False}{\bot}
\providecommand{\Args}[1]{#1}
\providecommand{\Ts}{\ensuremath{\top}}
\providecommand{\Fi}{\ensuremath{\bot}}
\providecommand{\lfalsum}{\bot}

% ---------------- objects, structures, metatheory ----------------
\providecommand{\Obj}{}
\providecommand{\Struct}{\ensuremath{\mathfrak{A}}}
\providecommand{\Domain}{\ensuremath{|\mathfrak{A}|}}
\providecommand{\Assign}{\ensuremath{s}}
\providecommand{\Sat}[2]{\ensuremath{#1 \models #2}}
\providecommand{\mSat}[3]{\ensuremath{#1 \models #2[#3]}}
\providecommand{\mModel}[2]{\ensuremath{#1 \models #2}}
\providecommand{\Lang}{\ensuremath{\mathcal{L}}}
\providecommand{\Th}{\ensuremath{\mathrm{Th}}}
\providecommand{\Log}{\ensuremath{\mathrm{L}}}
\providecommand{\tuple}[1]{\langle #1 \rangle}
\providecommand{\num}[1]{\underline{#1}}
\providecommand{\cfind}[1]{\varphi_{#1}}
\providecommand{\Nat}{\ensuremath{\mathbb{N}}}
\providecommand{\Mult}{\ensuremath{\mathrm{Mult}}}
\providecommand{\Add}{\ensuremath{\mathrm{Add}}}
\providecommand{\Exp}{\ensuremath{\mathrm{Exp}}}
\providecommand{\Zero}{\ensuremath{0}}
\providecommand{\Succ}{\ensuremath{\mathrm{S}}}
\providecommand{\lambd}{\ensuremath{\lambda}}
\providecommand{\bredpar}{\mathrel{\rightarrow_\beta}}
\providecommand{\bcd}[1]{\ensuremath{\ulcorner #1 \urcorner}}
\providecommand{\Set}[1]{\ensuremath{\{#1\}}}
\providecommand{\indcase}[1]{\par\noindent\emph{Case #1}.\quad}
\providecommand{\indfrm}[1]{\par\noindent\emph{Form #1}.\quad}
\providecommand{\iftag}[1]{#1}
\providecommand{\tagitem}[1]{\item[#1]}
\providecommand{\Elim}[1]{}
\providecommand{\Intro}[1]{}
\providecommand{\LeftR}[1]{}
\providecommand{\RightR}[1]{}
\providecommand{\TAss}[1]{}
\providecommand{\TRule}[1]{}

\makeatother
\providecommand{\olpath}{.}

% ===================== OLP compatibility layer (2) =====================
% --- second compatibility block: notation and shorthands used by the frozen OLP units ---
\providecommand{\CutCS}{\mathrm{Cut}}
\providecommand{\LeftSubproofLabel}{\relax}
\providecommand{\RightSubproofLabel}{\relax}
\providecommand{\maeh}{\relax}
\providecommand{\cheight}{\relax}
\providecommand{\ande}{\mathrm{and}_E}
\providecommand{\andi}{\mathrm{and}_I}
\providecommand{\ore}{\mathrm{or}_E}
\providecommand{\ori}{\mathrm{or}_I}
\providecommand{\typeof}{\mathord{:}}
\providecommand{\pValue}[1]{#1}
\providecommand{\Sent}{\mathrm{Sent}}
\providecommand{\cardplus}{\mathbin{+}}
\providecommand{\cardtimes}{\mathbin{\cdot}}
\providecommand{\cardexpo}{\mathbin{\hat{\ }}}
\providecommand{\cardnless}{\not<}
\providecommand{\cardle}{\leq}
\providecommand{\funfromto}[3]{#1\colon #2 \to #3}
\providecommand{\funimage}[1]{\operatorname{im}(#1)}
\providecommand{\fdefined}{\mathord{\downarrow}}
\providecommand{\canonord}{\operatorname{ord}}
\providecommand{\ordtimes}{\mathbin{\cdot}}
\providecommand{\ordeq}[2]{#1 = #2}
\providecommand{\ordneq}[2]{#1 \neq #2}
\providecommand{\ordtype}[1]{\operatorname{ord}(#1)}
\providecommand{\ordsucc}{\mathord{+1}}
\providecommand{\ordunion}{\cup}
\providecommand{\onesphere}{S^1}
\providecommand{\rotationsgroup}{\mathrm{SO}(2)}
\providecommand{\equivrep}[1]{[#1]}
\providecommand{\rlexless}{\prec}
\providecommand{\ZFC}{\mathrm{ZFC}}
\providecommand{\ZF}{\mathrm{ZF}}
\providecommand{\Rat}{\mathbb{Q}}
\providecommand{\Real}{\mathbb{R}}
\providecommand{\PosInt}{\mathbb{Z}^{+}}
\providecommand{\Id}{\mathrm{id}}
\providecommand{\Bin}{\mathrm{Bin}}
\providecommand{\pto}{\rightharpoonup}
\providecommand{\dom}{\operatorname{dom}}
\providecommand{\ran}{\operatorname{ran}}
\providecommand{\len}{\operatorname{len}}
\providecommand{\concat}{\mathbin{{}^\frown}}
\providecommand{\card}[1]{\lvert #1 \rvert}
\providecommand{\Pow}{\mathcal{P}}
\providecommand{\Setabs}[2]{\{\, #1 \mid #2 \,\}}
\providecommand{\comp}{\circ}
\providecommand{\Char}{\chi}
\providecommand{\Prf}{\operatorname{Prf}}
\providecommand{\Prov}{\operatorname{Prov}}
\providecommand{\Proj}{\operatorname{proj}}
\providecommand{\Gn}{\operatorname{gn}}
\providecommand{\gn}{\operatorname{gn}}
\providecommand{\Value}[1]{\operatorname{val}(#1)}
\providecommand{\Deduce}{\mathrm{D}}
\providecommand{\Dual}{\mathrm{d}}
\providecommand{\Ax}{\mathrm{Ax}}
\providecommand{\pSat}[3]{#1 \models #2\,[#3]}
\providecommand{\pAssign}{\sigma}
\providecommand{\bforall}{\forall}
\providecommand{\bexists}{\exists}
\providecommand{\liff}{\mathbin{\leftrightarrow}}
\providecommand{\FV}{\operatorname{FV}}
\providecommand{\UnaryInf}{\ensuremath{\mid}}
\providecommand{\BinaryInf}{\ensuremath{\mid}}
\providecommand{\depth}{\operatorname{depth}}
\providecommand{\oliflabeldef}[3]{#1}
\providecommand{\TMblank}{\sqcup}
\providecommand{\TMstroke}{\mid}
\providecommand{\TMstay}{\mathrm{S}}
\providecommand{\TMright}{\mathrm{R}}
\providecommand{\TMleft}{\mathrm{L}}
\providecommand{\TMendtape}{\triangleright}
\providecommand{\umin}{\mu}
\providecommand{\Sum}{\sum}
\providecommand{\citep}[1]{\cite{#1}}
\providecommand{\citet}[1]{\cite{#1}}
\providecommand{\citeauthor}[1]{\cite{#1}}
\providecommand{\citeyear}[1]{\cite{#1}}

% ===================== OLP compatibility layer (3) =====================
% --- third compatibility block: remaining notation seen in the frozen OLP units ---
\providecommand{\iso}{\cong}
\providecommand{\nszero}{\mathrm{0}}
\providecommand{\nssucc}{\mathrm{S}}
\providecommand{\nsplus}{\mathbin{+}}
\providecommand{\nssless}{<}
\providecommand{\nsless}{<}
\providecommand{\PL}{\mathrm{PL}}
\providecommand{\Nec}{\mathord{\Box}}
\providecommand{\RK}{\mathrm{RK}}
\providecommand{\ST}{\mathrm{ST}}
\providecommand{\SSubst}[3]{#1[#2/#3]}
\providecommand{\subst}[3]{#1[#2/#3]}
\providecommand{\tagfalse}{\mathrm{F}}
\providecommand{\cutr}{\mathrm{Cut}}
\providecommand{\cardneq}{\neq}
\providecommand{\issucc}[1]{\operatorname{succ}(#1)}
\providecommand{\tup}[1]{\langle #1 \rangle}
\providecommand{\Var}{\mathrm{Var}}
\providecommand{\Bew}{\operatorname{Bew}}
\providecommand{\Con}{\operatorname{Con}}
\providecommand{\Cons}{\operatorname{Cons}}
\providecommand{\pDeduce}{\mathrm{D}}
\providecommand{\psat}{\models}
\providecommand{\den}[1]{\llbracket #1 \rrbracket}
\providecommand{\Sem}[2]{\llbracket #1 \rrbracket_{#2}}
\providecommand{\DomainOf}[1]{\lvert #1 \rvert}
\providecommand{\Form}{\mathrm{Form}}
\providecommand{\Terms}{\mathrm{Term}}
\providecommand{\Atom}{\mathrm{Atom}}
\providecommand{\FVars}[1]{\operatorname{FV}(#1)}
\providecommand{\fv}[1]{\operatorname{fv}(#1)}
\providecommand{\tuplep}[1]{\langle #1 \rangle}
\providecommand{\restrict}[2]{#1 \upharpoonright #2}
\providecommand{\sto}{\rightarrow}
\providecommand{\FmlaSet}[1]{\mathrm{Fmla}_{#1}}
\providecommand{\Thm}{\mathrm{Thm}}
\providecommand{\Eval}{\operatorname{val}}
\providecommand{\blank}{\sqcup}
\providecommand{\blankTape}{\sqcup}
\providecommand{\Consis}{\operatorname{Con}}
\providecommand{\drs}{\mathrm{d}}
\providecommand{\VRule}{\mid}
\providecommand{\Bet}{\mathrm{B}}
\providecommand{\Seg}[1]{\mathrm{Seg}_{#1}}
\providecommand{\Enums}{\mathrm{En}}
\providecommand{\Z}{\mathbb{Z}}
\providecommand{\N}{\mathbb{N}}
\providecommand{\abs}[1]{\lvert #1 \rvert}
\providecommand{\floor}[1]{\lfloor #1 \rfloor}
\providecommand{\ceil}[1]{\lceil #1 \rceil}

% ===================== OLP compatibility layer (4) =====================
% --- fourth compatibility block: harvested from the per-candidate compile logs ---
\providecommand{\beredpar}{\mathrel{\rightarrow_\beta}}
\makeatletter
\def\bered{\@ifnextchar[{\bered@i}{\mathrel{\rightarrow_\beta}}}
\def\bered@i[#1]{\mathrel{\rightarrow_\beta}}
\def\bred{\@ifnextchar[{\bred@i}{\mathrel{\rightarrow_\beta}}}
\def\bred@i[#1]{\mathrel{\rightarrow_\beta}}
\makeatother
\providecommand{\beredone}{\mathrel{=_\beta}}
\providecommand{\bredone}{\mathrel{=_\beta}}
\providecommand{\xredone}{\mathrel{=_\beta}}
\providecommand{\aconvone}{\mathrel{=_\alpha}}
\providecommand{\cardless}{<}
\providecommand{\tagrefs}[1]{#1}
\providecommand{\ordplus}{\mathbin{+}}
\providecommand{\funrestrictionto}[2]{#1 \upharpoonright #2}
\providecommand{\fundefined}{\mathord{\uparrow}}
\providecommand{\Trm}{\mathrm{Trm}}
\providecommand{\OProv}{\operatorname{Prov}}
\providecommand{\OPrf}{\operatorname{Prf}}
\providecommand{\aeq}{\mathrel{=}}
\providecommand{\mClass}[1]{[#1]}
\providecommand{\equivclass}[1]{[#1]}
\providecommand{\pheight}{\operatorname{ht}}
\providecommand{\phy}{\operatorname{ht}}
\providecommand{\Cut}{\mathrm{Cut}}
\providecommand{\QuantRank}{\operatorname{qr}}
\providecommand{\elemequiv}{\mathrel{\equiv}}
\providecommand{\PIso}{\mathrm{PIso}}
\providecommand{\Theory}{\mathrm{Th}}
\providecommand{\supstrict}{\subsetneq}
\providecommand{\trcl}{\operatorname{trcl}}
\providecommand{\cref}[1]{\ref{#1}}
\providecommand{\Discharge}{\mathrm{Dis}}
\providecommand{\Weakening}{\mathrm{Wk}}
\providecommand{\ZFminus}{\mathrm{ZF}^{-}}
\makeatletter
\renewcommand{\Dual}{\@ifnextchar[{\Dual@i}{\ifmmode\mathrm{d}\else$\mathrm{d}$\fi}}
\def\Dual@i[#1]{\ifmmode\mathrm{d}[#1]\else$\mathrm{d}[#1]$\fi}
\makeatother
\providecommand{\RightR}{\mathrm{R}}
\providecommand{\LeftR}{\mathrm{L}}
\providecommand{\Intro}{\mathrm{I}}
\providecommand{\Axiom}{\mathrm{Ax}}
\providecommand{\formula}[1]{#1}
\providecommand{\stared}{\ast}
\providecommand{\oset}[2]{\{\, #1 \mid #2 \,\}}
\providecommand{\setof}[2]{\{\, #1 \mid #2 \,\}}
\providecommand{\pow}[1]{\mathcal{P}(#1)}
\providecommand{\ordpair}[2]{\langle #1, #2 \rangle}
\providecommand{\ordtriple}[3]{\langle #1, #2, #3 \rangle}
\providecommand{\union}{\cup}
\providecommand{\inter}{\cap}
\providecommand{\seq}[1]{\langle #1 \rangle}
\providecommand{\SuchThat}{\mid}
\providecommand{\where}{\mid}
\providecommand{\TVar}{\mathrm{TVar}}
\providecommand{\Bool}{\mathrm{Bool}}
\providecommand{\Val}{\operatorname{Val}}
\providecommand{\TL}[1]{\mathrm{TL}}
\providecommand{\xstared}{\ast}

% ===================== OLP compatibility layer (5) =====================
% --- fifth compatibility block: rescues the quarantined OLP candidates ---
\providecommand{\ident}{\operatorname{id}}
\providecommand{\setrank}{\operatorname{rank}}
\providecommand{\cardfont}{\mathfrak{c}}
\providecommand{\Complement}{\overline}
\providecommand{\QR}{\mathrm{QR}}
\providecommand{\becd}{\ensuremath{\ulcorner}}
\providecommand{\disjointsum}{\mathbin{\oplus}}
\providecommand{\cardeq}[2]{#1 = #2}
\providecommand{\xred}{\mathrel{\rightarrow}}
% --- sixth compatibility block: names reported by real compile logs ---
\providecommand{\ORefut}{\operatorname{Refut}}
\providecommand{\ORProv}{\operatorname{Prov}}
\providecommand{\cutrank}{\operatorname{cr}}
\providecommand{\Refut}{\operatorname{Refut}}
\providecommand{\DischargeRule}{\mathrm{Dis}}
\providecommand{\Contraction}{\mathrm{C}}
\providecommand{\bottomAlignProof}{\relax}
\providecommand{\FalseCl}{\bot}
\providecommand{\red}{\mathrel{\rightarrow}}
\providecommand{\varAssign}{\sigma}
\providecommand{\defiff}{\mathrel{\leftrightarrow}}
\providecommand{\FalseInt}{\bot}
\providecommand{\fn}{\operatorname{fn}}
\providecommand{\dotted}{\cdot}

% --- tag-macro semantic fix (2026-10-13, final): \iftag CHOOSES a branch and the OLP proof body is
% always the SECOND group (\iftag{tag}{body}, {tag}{body}{c}, {tag}{body}{c}{d}); a definition that
% prints the tag literal drops the proof body. Bodies contain blank lines, hence \long.
\makeatletter
\providecommand{\iftag}{}
\long\def\iftag#1#2{\@ifnextchar\bgroup{\iftag@three{#2}}{\iftag@stop{#2}}}
\long\def\iftag@three#1#2{\iftag@stop{#1}\@ifnextchar\bgroup{\iftag@eat}{\relax}}
\long\def\iftag@eat#1{}
\long\def\iftag@stop#1{#1}
\long\def\oliflabeldef#1#2{\@ifnextchar\bgroup{\@gobble}{\iftag@stop{#2}}}
\providecommand{\tagrefs}[1]{}
\providecommand{\settag}[1]{}
\makeatother
\begin{document}

% ==== OLP file: content/first-order-logic/syntax-and-semantics/assignments.tex ====
% Part: first-order-logic
% Chapter: syntax-and-semantics
% Section: assignments



\olfileid{fol}{syn}{ass}

\olsection{Variable Assignments}

\begin{explain}
A \olup{variable} assignment~$s$ provides a value for \emph{every}
variable---and there are infinitely many of them. This is of course
not necessary. We require \olup{variable} assignments to assign values to
all \olup{variable}s simply because it makes things a lot easier.  The
value of a term~$t$, and whether or not \olup{formula}~$A$ is
satisfied in \olup{structure} with respect to~$s$, only depend on the
assignments~$s$ makes to the \olup{variable}s in~$t$ and the free
\olup{variable}s of~$A$.  This is the content of the next two
propositions.  To make the idea of ``depends on'' precise, we show
that any two variable assignments that agree on all the variables
in~$t$ give the same value, and that $A$ is satisfied relative to one
iff it is satisfied relative to the other if two variable assignments
agree on all free variables of~$A$.
\end{explain}

\begin{prop}\ollabel{prop:valindep}
If the \olup{variable}s in a term~$t$ are among $x_1$, \dots,~$x_n$, and
$s_1(x_i) = s_2(x_i)$ for $i = 1$, \dots,~$n$, then $\Value{t}{M}[s_1]
= \Value{t}{M}[s_2]$.
\end{prop}

\begin{proof}
By induction on the complexity of~$t$. For the base case, $t$ can be
\olup{constant} or one of the variables~$x_1$, \dots,~$x_n$.  If $t
= c$, then $\Value{t}{M}[s_1] = \Assign{c}{M} = \Value{t}{M}[s_2]$. If
$t = x_i$, $s_1(x_i) = s_2(x_i)$ by the hypothesis of the proposition,
and so $\Value{t}{M}[s_1] = s_1(x_i) = s_2(x_i) = \Value{t}{M}[s_2]$.

For the inductive step, assume that $t = \Atom{f}{t_1, \dots, t_k}$
and that the claim holds for $t_1$, \dots, $t_k$. Then
\begin{align*}
  \Value{t}{M}[s_1] & = \Value{\Atom{f}{t_1, \dots, t_k}}{M}[s_1] \\
  & = \Assign{f}{M}(\Value{t_1}{M}[s_1], \dots, \Value{t_k}{M}[s_1]).
\intertext{For $j = 1$, \dots,~$k$, the \olup{variable}s of~$t_j$ are
    among $x_1$, \dots,~$x_n$. By the induction hypothesis,
    $\Value{t_j}{M}[s_1] = \Value{t_j}{M}[s_2]$. So,}
\Value{t}{M}[s_1] & = \Value{\Atom{f}{t_1, \dots, t_k}}{M}[s_1] \\
  & = \Assign{f}{M}(\Value{t_1}{M}[s_1], \dots, \Value{t_k}{M}[s_1]) \\
  & = \Assign{f}{M}(\Value{t_1}{M}[s_2], \dots, \Value{t_k}{M}[s_2]) \\
  & = \Value{\Atom{f}{t_1, \dots, t_k}}{M}[s_2] = \Value{t}{M}[s_2].
\end{align*}
\end{proof}

\begin{prop}\ollabel{prop:satindep}
If the free \olup{variable}s in $A$ are among $x_1$, \dots,~$x_n$, and
$s_1(x_i) = s_2(x_i)$ for $i = 1$, \dots,~$n$, then $\Sat{M}{A}[s_1]$
iff $\Sat{M}{A}[s_2]$.
\end{prop}

\begin{proof}
We use induction on the complexity of $A$. For the base case, where
$A$ is atomic, $A$ can be:
\iftag{prvTrue}{$\ltrue$,}{}
\iftag{prvFalse}{$\lfalse$,}{}
$\Atom{R}{t_1, \dots, t_k}$ for a $k$-place predicate $R$ and terms
$t_1$, \dots,~$t_k$, or $\eq[t_1][t_2]$ for terms $t_1$ and~$t_2$.
In the latter two cases, we only demonstrate the forward direction of
the \olup{biconditional}, since the proof of the reverse is symmetrical.

\begin{enumerate}
\tagitem{prvTrue}{%
  \indcase{A}{\ltrue}{both $\Sat{M}{\indfrm}[s_1]$ and
    $\Sat{M}{\indfrm}[s_2]$.}}{}

\tagitem{prvFalse}{%
  \indcase{A}{\lfalse}{both $\Sat/{M}{A}[s_1]$ and
    $\Sat/{M}{A}[s_2]$.}}{}

\item
  \indcase{A}{\Atom{R}{t_1, \ldots, t_k}}{let
    $\Sat{M}{\indfrm}[s_1]$. Then
    \[
    \langle \Value{t_1}{M}[s_1], \ldots, \Value{t_k}{M}[s_1] \rangle
    \in \Assign{R}{M}.
    \]
    For $i = 1$, \dots,~$k$, $\Value{t_i}{M}[s_1] =
    \Value{t_i}{M}[s_2]$ by \olref{prop:valindep}. So we also have
    $\langle \Value{t_i}{M}[s_2], \ldots, \Value{t_k}{M}[s_2] \rangle
    \in \Assign{R}{M}$, and hence $\Sat{M}{\indfrm}[s_2]$.}
\item
  \indcase{A}{\eq[t_1][t_2]}{suppose $\Sat{M}{\indfrm}[s_1]$.
    Then $\Value{t_1}{M}[s_1] = \Value{t_2}{M}[s_1]$. So,
    \begin{align*}
      \Value{t_1}{M}[s_2] & = \Value{t_1}{M}[s_1]
      & \text{(by \olref{prop:valindep})} \\
      & = \Value{t_2}{M}[s_1] 
      & \text{(since $\Sat{M}{\eq[t_1][t_2]}[s_1]$)}\\
      &= \Value{t_2}{M}[s_2]
      & \text{(by \olref{prop:valindep}),}
    \end{align*}
    so $\Sat{M}{\eq[t_1][t_2]}[s_2]$.}
\end{enumerate}

Now assume $\Sat{M}{B}[s_1]$ iff $\Sat{M}{B}[s_2]$ for all
\olup{formula}s $B$ less complex than~$A$. The induction step proceeds
by cases determined by the main operator of~$A$. In each case, we
only demonstrate the forward direction of the \olup{biconditional}; the
proof of the reverse direction is symmetrical. In all cases except
those for the quantifiers, we apply the induction hypothesis to
sub-\olup{formula}s~$B$ of~$A$.  The free variables of~$B$ are among
those of~$A$. Thus, if $s_1$ and $s_2$ agree on the free variables
of~$A$, they also agree on those of~$B$, and the induction
hypothesis applies to~$B$.

\begin{enumerate}
\tagitem{defNot}{}{%
  \iftag{probNot}{%
    \indcase!{A}{\lnot B}{}}{%
    \indcase{A}{\lnot B}{if $\Sat{M}{\indfrm}[s_1]$, then
      $\Sat/{M}{B}[s_1]$, so by the induction hypothesis,
      $\Sat/{M}{B}[s_2]$, hence $\Sat{M}{\indfrm}[s_2]$.}}}

\tagitem{defAnd}{}{%
  \iftag{probAnd}{%
    \indcase!{A}{B \land C}{}}{%
    \indcase{A}{B \land C}{if $\Sat{M}{\indfrm}[s_1]$, then
      $\Sat{M}{B}[s_1]$ and $\Sat{M}{C}[s_1]$, so by induction
      hypothesis, $\Sat{M}{B}[s_2]$ and $\Sat{M}{C}[s_2]$. Hence,
      $\Sat{M}{\indfrm}[s_2]$.}}}

\tagitem{defOr}{}{%
  \iftag{probOr}{%
    \indcase!{A}{B \lor C}{}}{%
    \indcase{A}{B \lor C}{if $\Sat{M}{\indfrm}[s_1]$, then
      $\Sat{M}{B}[s_1]$ or $\Sat{M}{C}[s_1]$. By induction hypothesis,
      $\Sat{M}{B}[s_2]$ or $\Sat{M}{C}[s_2]$, so $\Sat{M}{\indfrm}[s_2]$.}}}

\tagitem{defIf}{}{%
  \iftag{probIf}{%
    \indcase!{A}{B \lif C}{}}{%
    \indcase{A}{B \lif C}{if $\Sat{M}{\indfrm}[s_1]$, then
    $\Sat/{M}{B}[s_1]$ or $\Sat{M}{C}[s_1]$. By the induction hypothesis,
    $\Sat/{M}{B}[s_2]$ or $\Sat{M}{C}[s_2]$, so $\Sat{M}{\indfrm}[s_2]$.}}}

\tagitem{defIff}{}{%
  \iftag{probIff}{%
    \indcase!{A}{B \liff C}{}}{%
    \indcase{A}{B \liff C}{if $\Sat{M}{\indfrm}[s_1]$, then either
      $\Sat{M}{B}[s_1]$ and $\Sat{M}{C}[s_1]$, or $\Sat/{M}{B}[s_1]$ and
      $\Sat/{M}{C}[s_1]$. By the induction hypothesis, either
      $\Sat{M}{B}[s_2]$ and $\Sat{M}{C}[s_2]$ or $\Sat/{M}{B}[s_2]$
      and $\Sat/{M}{C}[s_2]$. In either case, $\Sat{M}{\indfrm}[s_2]$.}}}

\tagitem{defEx}{}{%
  \iftag{probEx}{%
    \indcase!{A}{\lexists[x][B]}{}}{%
    \indcase{A}{\lexists[x][B]}{if $\Sat{M}{\indfrm}[s_1]$, there is
      an $m \in \Domain{M}$ so that $\Sat{M}{B}[\Subst{s_1}{m}{x}]$.
      Let $s_1' = \Subst{s_1}{m}{x}$ and $s_2' =
      \Subst{s_2}{m}{x}$. The free variables of~$B$ are among $x_1$,
      \dots, $x_n$, and $x$. $s_1'(x_i) = s_2'(x_i)$, since $s_1'$ and
      $s_2'$ are $x$-variants of $s_1$ and~$s_2$, respectively, and by
      hypothesis $s_1(x_i) = s_2(x_i)$. $s_1'(x) = s_2'(x) = m$ by the
      way we have defined $s_1'$ and~$s_2'$. Then the induction
      hypothesis applies to $B$ and $s_1'$, $s_2'$, so
      $\Sat{M}{B}[s_2']$. Hence, since $s_2' = \Subst{s_2}{m}{x}$,
      there is an $m \in \Domain{M}$ such that
      $\Sat{M}{B}[\Subst{s_2}{m}{x}]$, and so
      $\Sat{M}{\indfrm}[s_2]$.}}}

\tagitem{defAll}{}{%
  \iftag{probAll}{%
    \indcase!{A}{\lforall[x][B]}{}}{%
    \indcase{A}{\lforall[x][B]}{if $\Sat{M}{\indfrm}[s_1]$, then for
      every $m \in \Domain{M}$, $\Sat{M}{B}[\Subst{s_1}{m}{x}]$. We
      want to show that also, for every $m \in \Domain{M}$,
      $\Sat{M}{B}[\Subst{s_2}{m}{x}]$. So let $m \in \Domain{M}$ be
      arbitrary, and consider $s_1' = \Subst{s}{m}{x}$ and $s_2' =
      \Subst{s}{m}{x}$. We have that $\Sat{M}{B}[s_1']$. The free
      variables of~$B$ are among $x_1$, \dots, $x_n$, and $x$.
      $s_1'(x_i) = s_2'(x_i)$, since $s_1'$ and $s_2'$ are
      $x$-variants of $s_1$ and~$s_2$, respectively, and by hypothesis
      $s_1(x_i) = s_2(x_i)$. $s_1'(x) = s_2'(x) = m$ by the way we
      have defined $s_1'$ and~$s_2'$. Then the induction hypothesis
      applies to~$B$ and~$s_1'$, $s_2'$, and we have
      $\Sat{M}{B}[s_2']$. This applies to every $m \in \Domain{M}$,
      i.e., $\Sat{M}{B}[\Subst{s_2}{m}{x}]$ for all~$m \in
      \Domain{M}$, so $\Sat{M}{\indfrm}[s_2]$.}}}
\end{enumerate}
By induction, we get that $\Sat{M}{A}[s_1]$ iff $\Sat{M}{A}[s_2]$
whenever the free \olup{variable}s in $A$ are among $x_1$, \dots, $x_n$
and $s_1(x_i)=s_2(x_i)$ for $i = 1$, \dots,~$n$.
\end{proof}

\begin{probtag}{probNot,probOr,probAnd,probIf,probIff,probEx,probAll}
Complete the proof of \olref[fol][syn][ass]{prop:satindep}.
\end{probtag}

\begin{explain}
!!^{sentence}s have no free variables, so any two variable assignments  
assign the same things to all the (zero) free variables of any
sentence. The proposition just proved then means that whether or not
\olup{sentence} is satisfied in a structure relative to a variable
assignment is completely independent of the assignment. We'll record
this fact. It justifies the definition of satisfaction of
\olup{sentence} in \olup{structure} (without mentioning a variable
assignment) that follows.
\end{explain}

\begin{cor}
\ollabel{cor:sat-sentence}
If $A$ is \olup{sentence} and $s$ a variable assignment, then
$\Sat{M}{A}[s]$ iff $\Sat{M}{A}[s']$ for every variable
assignment~$s'$.
\end{cor}

\begin{proof}
  Let $s'$ be any variable assignment. Since $A$ is \olup{sentence}, it
  has no free variables, and so every variable assignment~$s'$
  trivially assigns the same things to all free variables of~$A$ as
  does~$s$. So the condition of \olref{prop:satindep} is satisfied,
  and we have $\Sat{M}{A}[s]$ iff $\Sat{M}{A}[s']$.
\end{proof}

\begin{defn}
\ollabel{defn:satisfaction}
If $A$ is \olup{sentence}, we say that \olup{structure}~$\Struct M$
\emph{satisfies}~$A$, $\Sat{M}{A}$, iff $\Sat{M}{A}[s]$ for all
variable assignments~$s$.
\end{defn}

If $\Sat{M}{A}$, we also simply say that \emph{$A$ is true
in~$\Struct{M}$.} The notion of satisfaction naturally extends
from individual \olup{sentence}s to sets of \olup{sentence}s.

\begin{defn}
\ollabel{defn:sat}
If $\Gamma$ is a set of \olup{sentence}s~$\Gamma$, we say that
\olup{structure}~$\Struct M$ \emph{satisfies}~$\Gamma$,
$\Sat{M}{\Gamma}$, iff $\Sat{M}{A}$ for all $A \in \Gamma$.
\end{defn}

\begin{prop}\ollabel{prop:sentence-sat-true}
  Let $\Struct{M}$ be \olup{structure}, $A$ be \olup{sentence}, and $s$ a
  variable assignment.  $\Sat{M}{A}$ iff $\Sat{M}{A}[s]$.
\end{prop}

\begin{proof}
Exercise.
\end{proof}

\begin{prob}
Prove \olref[fol][syn][ass]{prop:sentence-sat-true}
\end{prob}

\begin{prop}\ollabel{prop:sat-quant}
  Suppose $A(x)$ only contains $x$ free, and $\Struct M$ is
  \olup{structure}. Then:
  \begin{tagenumerate}{prvEx,prvAll}
    \tagitem{prvEx}{$\Sat{M}{\lexists[x][A(x)]}$ iff $\Sat{M}{A(x)}[s]$
      for at least one variable assignment~$s$.}{}
    \tagitem{prvAll}{$\Sat{M}{\lforall[x][A(x)]}$ iff $\Sat{M}{A(x)}[s]$
  for all variable assignments~$s$.}{}
\end{tagenumerate}
\end{prop}

\begin{proof}
Exercise.
\end{proof}

\begin{prob}
Prove \olref[fol][syn][ass]{prop:sat-quant}.
\end{prob}

\begin{prob}
\DeclareRobustCommand{\VDash}{\mathrel{||}\joinrel\Relbar}  
Suppose $\Lang L$ is a language without \olup{function}s. Given a
\olup{structure}~$\Struct M$, $c$ \olup{constant} and $a \in \Domain M$,
define $\Struct M[a/c]$ to be the \olup{structure} that is just
like~$\Struct M$, except that $\Assign{c}{M[a/c]} = a$. Define
$\Struct M \VDash A$ for \olup{sentence}s~$A$ by:
\begin{enumerate}
\tagitem{prvFalse}{%
  \indcase{A}{\lfalse}{not $\Struct{M} \VDash \indfrm$.}}{}

\tagitem{prvTrue}{%
  \indcase{A}{\ltrue}{$\Struct{M} \VDash \indfrm$.}}{}

\item \indcase{A}{\Atom{R}{d_1, \dots, d_n}}{$\Struct M \VDash \indfrm$
  iff $\langle \Assign{d_1}{M}, \dots, \Assign{d_n}{M} \rangle \in
  \Assign{R}{M}$.}

\item \indcase{A}{\eq[d_1][d_2]}{$\Struct M \VDash \indfrm$ iff
  $\Assign{d_1}{M} = \Assign{d_2}{M}$.}

\tagitem{prvNot}{%
  \indcase{A}{\lnot B}{$\Struct{M} \VDash \indfrm$ iff
    not $\Struct{M} \VDash {B}$.}}{}

\tagitem{prvAnd}{%
  \indcase{A}{(B \land C)}{$\Struct{M} \VDash
    \indfrm$ iff $\Struct{M} \VDashB$ and $\Struct{M} \VDash C$.}}{}

\tagitem{prvOr}{%
  \indcase{A}{(B \lor C)}{$\Struct{M} \VDash \indfrm$ iff
    $\Struct{M} \VDash B$ or $\Struct{M} \VDash C$ (or both).}}{}

\tagitem{prvIf}{%
  \indcase{A}{(B \lif C)}{$\Struct{M} \VDash \indfrm$ iff
    not $\Struct {M} \VDash B$ or $\Struct M \VDash C$ (or both).}}{}

\tagitem{prvIff}{%
  \indcase{A}{(B \liff C)}{$\Struct{M} \VDash {\indfrm}$ iff either
    both $\Struct{M} \VDash {B}$ and $\Struct{M} \VDash {C}$, or
    neither $\Struct{M} \VDash {B}$ nor $\Struct{M} \VDash {C}$.}}{}

\tagitem{prvAll}{%
  \indcase{A}{\lforall[x][B]}{$\Struct{M} \VDash {\indfrm}$ iff for
    all $a \in \Domain{M}$, $\Struct{M[a/c]} \VDash \Subst{B}{c}{x}$,
    if $c$ does not occur in~$B$.}}{}

\tagitem{prvEx}{%
  \indcase{A}{\lexists[x][B]}{$\Struct{M} \VDash
  {\indfrm}$ iff there is an $a \in \Domain M$ such that
  $\Struct{M[a/c]} \VDash \Subst{B}{c}{x}$, if $c$ does not occur
  in~$B$.}}{}
\end{enumerate}
Let $x_1$, \dots, $x_n$ be all free \olup{variable}s in~$A$,
$c_1$, \dots, $c_n$ constant symbols not in~$A$,
$a_1$, \dots, $a_n \in \Domain M$, and $s(x_i) = a_i$.

Show that $\Sat{M}{A}[s]$ iff $\Struct M[a_1/c_1,\dots,a_n/c_n]
\VDash \Subst{\Subst{A}{c_1}{x_1}\dots}{c_n}{x_n}$.

(This problem shows that it is possible to give a semantics for
first-order logic that makes do without variable assignments.)
\end{prob}

\begin{prob}
Suppose that $f$ is a function symbol not in~$A(x,y)$. Show that
there is \olup{structure}~$\Struct{M}$ such that
$\Sat{M}{\lforall[x][\lexists[y][A(x,y)]]}$ iff there is an~$\Struct
M'$ such that $\Sat{M'}{\lforall[x][A(x,f(x))]}$.

(This problem is a special case of what's known as Skolem's Theorem;
$\lforall[x][A(x,f(x))]$ is called a \emph{Skolem normal form} of
$\lforall[x][\lexists[y][A(x,y)]]$.)
\end{prob}


% ==== OLP file: content/first-order-logic/syntax-and-semantics/covered-structures.tex ====
% Part: first-order-logic
% Chapter: syntax-and-semantics
% Section: covered-structures



\olfileid{fol}{syn}{cov}

\olsection{Covered \printtoken{P}{structure} for First-order Languages}

\begin{explain}
Recall that a term is \emph{closed} if it contains no \olup{variable}s.
\end{explain}

\begin{defn}[!!^{value} of closed terms]
If $t$ is a closed term of the language~$\Lang L$ and $\Struct M$ is a
\olup{structure} for~$\Lang L$, the \emph{\olup{value}}~$\Value{t}{M}$ is
defined as follows:
\begin{enumerate}
\item If $t$ is just the \olup{constant}~$c$, then $\Value{c}{M} = \Assign{c}{M}$.
\item If $t$ is of the form $\Atom{f}{t_1, \ldots, t_n}$, then
  \[
  \Value{t}{M} = \Assign{f}{M}(\Value{t_1}{M}, \ldots,
  \Value{t_n}{M}).
  \]
\end{enumerate}
\end{defn}

\begin{defn}[Covered \olup{structure}]
A \olup{structure} is \emph{covered} if every element of the domain is the
\olup{value} of some closed term.
\end{defn}

\begin{ex}
Let ~$\Lang L$ be the language with \olup{constant}s $\Obj{zero}$,
$\Obj{one}$, $\Obj{two}$, \dots, the binary \olup{predicate}~$<$, and the
binary \olup{function}s $+$ and $\times$. Then \olup{structure}~$\Struct
M$ for~$\Lang L$ is the one with domain $\Domain M = \{0, 1, 2, \ldots
\}$ and assignments $\Assign{\Obj{zero}}{M} = 0$,
$\Assign{\Obj{one}}{M} = 1$, $\Assign{\Obj{two}}{M} = 2$, and so
forth. For the binary relation symbol $<$, the set $\Assign{<}{M}$ is
the set of all pairs $\tuple{c_1, c_2} \in \Domain{M}^2$ such that
$c_1$ is less than~$c_2$: for example, $\tuple{1, 3} \in
\Assign{<}{M}$ but $\tuple{2, 2} \notin \Assign{<}{M}$. For the binary
\olup{function} $+$, define $\Assign{+}{M}$ in the usual way---for
example, $\Assign{+}{M}(2,3)$ maps to~$5$, and similarly for the
binary \olup{function}~$\times$. Hence, the \olup{value} of $\Obj{four}$ is
just~$4$, and the \olup{value} of $\times(\Obj{two},
+(\Obj{three},\Obj{zero}))$ (or in infix notation, $\Obj{two} \times
(\Obj{three} + \Obj{zero})$) is
\begin{multline*}
\Value{\times(\Obj{two}, +(\Obj{three},\Obj{zero}))}{M} =\\
\begin{aligned}
& =\Assign{\times}{M}(\Value{\Obj{two}}{M}, \Value{+(\Obj{three}, \Obj{zero})}{M})\\
& = \Assign{\times}{M}(\Value{\Obj{two}}{M}, \Assign{+}{M}(\Value{\Obj{three}}{M},
\Value{\Obj{zero}}{M})) \\
& = \Assign{\times}{M}(\Assign{\Obj{two}}{M}, \Assign{+}{M}(\Assign{\Obj{three}}{M},
\Assign{\Obj{zero}}{M})) \\
& = \Assign{\times}{M}(2, \Assign{+}{M}(3, 0)) \\
& = \Assign{\times}{M}(2, 3) \\
& = 6
\end{aligned}
\end{multline*}
\end{ex}

\begin{prob}
Is $\Struct N$, the standard model of arithmetic, covered? Explain.
\end{prob}


% ==== OLP file: content/first-order-logic/syntax-and-semantics/extensionality.tex ====
% Part: first-order-logic
% Chapter: syntax-and-semantics
% Section: substitution



\olfileid{fol}{syn}{ext}

\olsection{Extensionality}

\begin{explain}
Extensionality, sometimes called relevance, can be expressed
informally as follows: the only factors that bear upon the
satisfaction of \olup{formula}~$A$ in \olup{structure}~$\Struct M$
relative to \olup{variable} assignment~$s$, are the size of the
\olup{domain} and the assignments made by~$\Struct M$ and~$s$ to the
elements of the language that actually appear in~$A$.

One immediate consequence of extensionality is that where two
\olup{structure}s~$\Struct M$ and~$\Struct M'$ agree on all the elements
of the language appearing in a sentence~$A$ and have the same
domain,~$\Struct M$ and~$\Struct M'$ must also agree on whether or not
$A$ itself is true.
\end{explain}

\begin{prop}[Extensionality]
  \ollabel{prop:extensionality}
  Let $A$ be \olup{formula}, and $\Struct M_1$ and $\Struct M_2$ be
  \olup{structure}s with $\Domain{M_1} = \Domain{M_2}$, and $s$ a
  variable assignment on $\Domain{M_1} = \Domain{M_2}$.  If
  $\Assign{c}{M_1} = \Assign{c}{M_2}$, $\Assign{R}{M_1}=\Assign{R}{M_2}$,
  and $\Assign{f}{M_1} = \Assign{f}{M_2}$ for every \olup{constant}~$c$,
  relation symbol~$R$, and \olup{function} $f$ occurring in~$A$, then
  $\Sat{M_1}{A}[s]$ iff $\Sat{M_2}{A}[s]$.
\end{prop}

\begin{proof}
  First prove (by induction on~$t$) that for every term,
  $\Value{t}{M_1}[s] = \Value{t}{M_2}[s]$.  Then prove the proposition
  by induction on~$A$, making use of the claim just proved for the
  induction basis (where $A$ is atomic).
\end{proof}

\begin{prob}
Carry out the proof of \olref[fol][syn][ext]{prop:extensionality} in
detail.
\end{prob}

\begin{cor}[Extensionality for !!^{sentence}s]
  \ollabel{cor:extensionality-sent}
  Let $A$ be \olup{sentence} and $\Struct{M_1}$, $\Struct{M_2}$ as in
  \olref{prop:extensionality}. Then $\Sat{M_1}{A}$ iff $\Sat{M_2}{A}$.
\end{cor}

\begin{proof}
Follows from \olref{prop:extensionality} by \olref[ass]{cor:sat-sentence}.
\end{proof}

Moreover, the value of a term, and whether or not \olup{structure}
satisfies \olup{formula}, only depend on the values of its subterms.

\begin{prop}\ollabel{prop:ext-terms}
Let $\Struct M$ be \olup{structure}, $t$ and $t'$ terms, and $s$ a
variable assignment. Then $\Value{\Subst{t}{t'}{x}}{M}[s] =
\Value{t}{M}[\Subst{s}{\Value{t'}{M}[s]}{x}]$.
\end{prop}

\begin{proof}
By induction on~$t$.
\begin{enumerate}
\item If $t$ is a constant, say, $t\ident c$, then $\Subst{t}{t'}{x} =
  c$, and $\Value{c}{M}[s] = \Assign{c}{M} =
  \Value{c}{M}[\Subst{s}{\Value{t'}{M}[s]}{x}]$.

\item If $t$ is a variable other than~$x$, say, $t \ident y$, then
  $\Subst{t}{t'}{x} = y$, and $\Value{y}{M}[s] =
  \Value{y}{M}[\Subst{s}{\Value{t'}{M}[s]}{x}]$ since
  $\varAssign{s}{\Subst{s}{\Value{t'}{M}[s]}{x}}{x}$.

\item If $t \ident x$, then $\Subst{t}{t'}{x} = t'$. But
  $\Value{x}{M}[\Subst{s}{\Value{t'}{M}[s]}{x}] = \Value{t'}{M}[s]$ by
  definition of~$\Subst{s}{\Value{t'}{M}[s]}{x}$.

\item If $t \ident \Atom{f}{t_1,\dots,t_n}$ then we have:
\begin{multline*}
  \Value{\Subst{t}{t'}{x}}{M}[s]  = \\
  \begin{aligned}[b]
& = \Value{\Atom{f}{\Subst{t_1}{t'}{x}, \dots, \Subst{t_n}{t'}{x}}}{M}[s]\\
& \qquad    \text{ by definition of $\Subst{t}{t'}{x}$}\\
& = \Assign{f}{M}(\Value{\Subst{t_1}{t'}{x}}{M}[s], \dots,
    \Value{\Subst{t_n}{t'}{x}}{M}[s])\\
    & \qquad  \text{ by definition of $\Value{\Atom{f}{\dots}}{M}[s]$}\\
& = \Assign{f}{M}(\Value{t_1}{M}[\Subst{s}{\Value{t'}{M}[s]}{x}], \dots,
   \Value{t_n}{M}[\Subst{s}{\Value{t'}{M}[s]}{x}])\\
& \qquad    \text{ by induction hypothesis}\\
& = \Value{t}{M}[\Subst{s}{\Value{t'}{M}[s]}{x}]
    \text{ by definition of $\Value{\Atom{f}{\dots}}{M}[\Subst{s}{\Value{t'}{M}[s]}{x}]$}
  \end{aligned}
\end{multline*}
\end{enumerate}
\end{proof}

\begin{prop}\ollabel{prop:ext-formulas} Let $\Struct M$ be
\olup{structure}, $A$ \olup{formula}, $t'$~a term, and $s$~a variable
assignment. Then $\Sat{M}{\Subst{A}{t'}{x}}[s]$ iff
$\Sat{M}{A}[\Subst{s}{\Value{t'}{M}[s]}{x}]$.
\end{prop}

\begin{proof}
Exercise.
\end{proof}

\begin{prob}
Prove \olref[fol][syn][ext]{prop:ext-formulas}
\end{prob}

\begin{explain}
  The point of
  \cref{fol:syn:ext:prop:ext-terms,fol:syn:ext:prop:ext-formulas} is
  the following. Suppose we have a term $t$ or \olup{formula}~$A$ and
  some term~$t'$, and we want to know the value of $\Subst{t}{t'}{x}$
  or whether or not $\Subst{A}{t'}{x}$ is satisfied in
  \olup{structure}~$\Struct M$ relative to \olup{variable} assignment~$s$.
  Then we can either perform the substitution first and then consider
  the value or satisfaction relative to $\Struct{M}$ and~$s$, or we
  can first determine the value~$m = \Value{t'}{M}[s]$ of $t'$ in
  $\Struct{M}$ relative to~$s$, change the \olup{variable} assignment
  to~$\Subst{s}{m}{x}$ and then consider the value of~$t$ in
  $\Struct{M}$ and~$\Subst{s}{m}{x}$, or whether
  $\Sat{M}{A}[\Subst{s}{m}{x}]$.
  \Cref{fol:syn:ext:prop:ext-terms,fol:syn:ext:prop:ext-formulas}
  guarantee that the answer will be the same, whichever way we do it.
\end{explain}


% ==== OLP file: content/first-order-logic/syntax-and-semantics/first-order-languages.tex ====
% Part: first-order-logic
% Chapter: syntax-and-semantics
% Section: first-order-languages



\olfileid{fol}{syn}{fol}

\olsection{First-Order Languages}


Expressions of first-order logic are built up from a basic vocabulary
containing \emph{\olup{variable}s}, \emph{\olup{constant}s},
\emph{\olup{predicate}s} and sometimes \emph{\olup{function}s}.  From them,
together with logical connectives, quantifiers, and punctuation
symbols such as parentheses and commas, \emph{terms} and
\emph{\olup{formula}s} are formed.

\begin{explain}
Informally, \olup{predicate}s are names for properties and relations,
\olup{constant}s are names for individual objects, and \olup{function}s are
names for mappings.  These, except for the \olup{identity}~$\eq$, are the
\emph{non-logical symbols} and together make up a language.  Any
first-order language~$\Lang L$ is determined by its non-logical
symbols.  In the most general case, $\Lang L$ contains infinitely many
symbols of each kind.
\end{explain}

In the general case, we make use of the following symbols in
first-order logic:

\begin{enumerate}
\item Logical symbols
\begin{enumerate}
\item Logical connectives:
  \startycommalist
  \iftag{prvNot}{\ycomma $\lnot$ (negation)}{}%
  \iftag{prvAnd}{\ycomma $\land$ (conjunction)}{}%
  \iftag{prvOr}{\ycomma $\lor$ (disjunction)}{}%
  \iftag{prvIf}{\ycomma $\lif$ (\olup{conditional})}{}%
  \iftag{prvIff}{\ycomma $\liff$ (\olup{biconditional})}{}%
  \iftag{prvAll}{\ycomma $\lforall$ (universal quantifier)}{}%
  \iftag{prvEx}{\ycomma $\lexists$ (existential quantifier)}{}.
\tagitem{prvFalse}{The propositional constant for \olup{falsity}~$\lfalse$.}{}
\tagitem{prvTrue}{The propositional constant for \olup{truth}~$\ltrue$.}{}
\item The two-place \olup{identity}~$\eq$.
\item A \olup{denumerable}s set of \olup{variable}s: $\Obj v_0$, $\Obj v_1$, $\Obj
  v_2$, \dots
\end{enumerate}
\item Non-logical symbols, making up the \emph{standard
  language} of first-order logic
\begin{enumerate}
\item A \olup{denumerable}s set of $n$-place \olup{predicate}s for each $n>0$: $\Obj
  A^n_0$, $\Obj A^n_1$, $\Obj A^n_2$, \dots
\item A \olup{denumerable}s set of \olup{constant}s: $\Obj c_0$, $\Obj c_1$, $\Obj
  c_2$, \dots.
\item A \olup{denumerable}s set of $n$-place \olup{function}s for each $n>0$:
  $\Obj f^n_0$, $\Obj f^n_1$, $\Obj f^n_2$, \dots
\end{enumerate}
\item Punctuation marks: (, ), and the comma.
\end{enumerate}

Most of our definitions and results will be formulated for the full
standard language of first-order logic.  However, depending on the
application, we may also restrict the language to only a few
\olup{predicate}s, \olup{constant}s, and \olup{function}s.

\begin{ex}
The language~$\Lang L_A$ of arithmetic contains a single two-place
\olup{predicate}~$<$, a single \olup{constant}~$\Obj 0$, one one-place
\olup{function}~$\prime$, and two two-place \olup{function}s~$+$ and~$\times$.
\end{ex}

\begin{ex}
The language of set theory~$\Lang L_Z$ contains only the single
two-place \olup{predicate}~$\in$.
\end{ex}

\begin{ex}
The language of orders~$\Lang L_\le$ contains only the two-place
\olup{predicate}~$\le$.
\end{ex}

Again, these are conventions: officially, these are just aliases,
e.g., $<$, $\in$, and $\le$ are aliases for $\Obj A^2_0$, $\Obj 0$ for
$\Obj c_0$, $\prime$ for $\Obj f^1_0$, $+$ for $\Obj f^2_0$, $\times$ for
$\Obj f^2_1$.

\iftag{defNot,defOr,defAnd,defIf,defIff,defTrue,defFalse,defEx,defAll}{%
In addition to the primitive connectives and
\iftag{notprvEx,notprvAll}{quantifier}{quantifiers} introduced
above, we also use the following \emph{defined} symbols:
\startycommalist
  \iftag{defNot}{\ycomma $\lnot$ (negation)}{}%
  \iftag{defAnd}{\ycomma $\land$ (conjunction)}{}%
  \iftag{defOr}{\ycomma $\lor$ (disjunction)}{}%
  \iftag{defIf}{\ycomma $\lif$ (\olup{conditional})}{}%
  \iftag{defIff}{\ycomma $\liff$ (\olup{biconditional})}{}%
  \iftag{defAll}{\ycomma $\lforall$ (universal quantifier)}{}%
  \iftag{defEx}{\ycomma $\lexists$ (existential quantifier)}{}%
  \iftag{defFalse}{\ycomma \olup{falsity}~$\lfalse$}{}%
  \iftag{defTrue}{\ycomma \olup{truth}~$\ltrue$}}{}.

\begin{tagblock}{defNot,defOr,defAnd,defIf,defIff,defTrue,defFalse,defEx,defAll}
\begin{explain}
A defined symbol is not officially part of the language, but is
introduced as an informal abbreviation: it allows us to abbreviate
formulas which would, if we only used primitive symbols, get quite
long.  This is obviously an advantage.  The bigger advantage, however,
is that proofs become shorter.  If a symbol is primitive, it has to be
treated separately in proofs. The more primitive symbols, therefore,
the longer our proofs.
\end{explain}
\end{tagblock}

% Alternate symbols

\begin{intro}
You may be familiar with different terminology and symbols than the
ones we use above. Logic texts (and teachers) commonly use 
$\sim$, $\neg$, or~!\ for ``negation'', $\wedge$, $\cdot$, or $\&$
for ``conjunction''.  Commonly used symbols for the ``conditional'' or
``implication'' are $\rightarrow$, $\Rightarrow$, and $\supset$.
\iftag{prvIff,defIff}{Symbols for ``biconditional,'' ``bi-implication,''
  or ``(material) equivalence'' are $\leftrightarrow$,
  $\Leftrightarrow$, and $\equiv$.}{}
\iftag{prvFalse,defFalse}{The $\lfalse$ symbol is variously called
  ``falsity,'' ``falsum,'', ``absurdity,'' or ``bottom.''}{}
\iftag{prvTrue,defTrue}{The $\ltrue$ symbol is variously called
  ``truth,'' ``verum,'' or ``top.''}{}

It is conventional to use lower case letters (e.g., $a$, $b$, $c$) from
the beginning of the Latin alphabet for \olup{constant}s (sometimes called
names), and lower case letters from the end (e.g., $x$, $y$, $z$) for
\olup{variable}s. Quantifiers combine with \olup{variable}s, e.g., $x$;
notational variations include $\forall x$, $(\forall x)$, $(x)$, $\Pi x$,
$\bigwedge_x$ for the universal quantifier and $\exists x$, $(\exists
x)$, $(Ex)$, $\Sigma x$, $\bigvee_x$ for the existential quantifier.
\end{intro}

\begin{explain}
We might treat all the propositional operators and both quantifiers as
primitive symbols of the language.  We might instead choose a smaller
stock of primitive symbols and treat the other \olup{operator}s as
defined. ``Truth functionally complete'' sets of Boolean operators
include $\{ \lnot, \lor \}$, $\{ \lnot, \land \}$, and $\{ \lnot,
\lif\}$---these can be combined with either quantifier for an
expressively complete first-order language.

You may be familiar with two other \olup{operator}s: the Sheffer
stroke~$|$ (named after Henry Sheffer), and Peirce's
arrow~$\downarrow$, also known as Quine's dagger.  When given their
usual readings of ``nand'' and ``nor'' (respectively), these operators
are truth functionally complete by themselves.
\end{explain}


% ==== OLP file: content/first-order-logic/syntax-and-semantics/formation-sequences.tex ====
% Part: first-order-logic
% Chapter: syntax-and-semantics
% Section: formation-sequences



\olfileid{fol}{syn}{fseq}

\olsection{Formation Sequences}

Defining \olup{formula}s via an inductive definition, and the
complementary technique of proving properties of \olup{formula}s via
induction, is an elegant and efficient approach. However, it can
also be useful to consider a more bottom-up, step-by-step approach
to the construction of \olup{formula}s, which we do here using the
notion of a \emph{formation sequence}.
%
To show how terms and \olup{formula}s can be introduced in this way
without needing to refer to their inductive definitions, we first
introduce the notion of an arbitrary string of symbols drawn from
some language~$\Lang L$.

\begin{defn}[Strings]
\ollabel{defn:string}
Suppose $\Lang L$ is a first-order language. An \emph{$\Lang
L$-string} is a finite sequence of symbols of~$\Lang L$. Where the
language~$\Lang L$ is clearly fixed by the context, we will often
refer to a $\Lang L$-string simply as a \emph{string}.
\end{defn}

\begin{ex}
For any first-order language $\Lang L$, all
$\Lang L$-\olup{formula}s are $\Lang L$-strings, but not
conversely. For example, \[)(\Obj v_0\lif\lexists\] is an
$\Lang L$-string but not an $\Lang L$-\olup{formula}.
\end{ex}

\begin{defn}[Formation sequences for terms]
\ollabel{defn:fseq-trm}
A finite sequence of $\Lang L$-strings $\tuple{t_0,\dotsc,t_n}$ is a
\emph{formation sequence} for a term $t$ if $t \ident t_n$ and for all
$i \leq n$, either $t_i$ is \olup{variable} or \olup{constant}, or $\Lang
L$ contains a $k$-ary \olup{function}~$f$ and there exist
$m_0,\dotsc,m_k < i$ such that $t_i \ident f(t_{m_0},\dotsc,t_{m_k})$.
When it is necessary to distinguish, we will refer to formation
sequences for terms as \emph{term formation sequences}.
\end{defn}

\begin{ex}
The sequence
\[
    \tuple{\Obj c_0, \Obj v_0, \Atom{\Obj f^2_0}{\Obj c_0, \Obj v_0}, \Atom{\Obj f^1_0}{\Atom{\Obj f^2_0}{\Obj c_0, \Obj v_0}}}
\]
is a formation sequence for the term $\Atom{\Obj f^1_0}{\Atom{\Obj
f^2_0}{\Obj c_0, \Obj v_0}}$, as is
\[
    \tuple{\Obj v_0, \Obj c_0, \Atom{\Obj f^2_0}{\Obj c_0, \Obj v_0}, \Atom{\Obj f^1_0}{\Atom{\Obj f^2_0}{\Obj c_0, \Obj v_0}}}.
\]
\end{ex}

\begin{defn}[Formation sequences for formulas]
\ollabel{defn:fseq-frm}
A finite sequence of $\Lang L$-strings $\tuple{A_0,\dotsc,A_n}$
is a \emph{formation sequence} for~$A$ if $A \ident A_n$ and
for all $i \leq n$, either $A_i$ is an atomic \olup{formula} or there
exist $j,k < i$ and \olup{variable}~$x$ such that one of the following
holds:
\begin{enumerate}
    \tagitem{prvNot}{$A_i \ident \lnot A_j$.}{}%
    \tagitem{prvAnd}{$A_i \ident (A_j \land A_k)$.}{}%
    \tagitem{prvOr}{$A_i \ident (A_j \lor A_k)$.}{}%
    \tagitem{prvIf}{$A_i \ident (A_j \lif A_k)$.}{}%
    \tagitem{prvIff}{$A_i \ident (A_j \liff A_k)$.}{}%
    \tagitem{prvAll}{$A_i \ident \lforall[x][A_j]$.}{}%
    \tagitem{prvEx}{$A_i \ident \lexists[x][A_j]$.}{}%
\end{enumerate}
When it is necessary to distinguish, we will refer to formation
sequences for formulas as \emph{formula formation sequences}.
\end{defn}

\begin{ex}
\[
\tuple{
    \Atom{\Obj A^1_0}{\Obj v_0},
    \Atom{\Obj A^1_1}{\Obj c_1},
    (\Atom{\Obj A^1_1}{\Obj c_1} \land \Atom{\Obj A^1_0}{\Obj v_0}),
    \lexists[\Obj v_0][(\Atom{\Obj A^1_1}{\Obj c_1} \land \Atom{\Obj A^1_0}{\Obj v_0})]
}
\]
is a formation sequence of $\lexists[\Obj v_0][(\Atom{\Obj A^1_1}{\Obj
c_1} \land \Atom{\Obj A^1_0}{\Obj v_0})]$, as is
\begin{multline*}
\tuple{
    \Atom{\Obj A^1_0}{\Obj v_0},
    \Atom{\Obj A^1_1}{\Obj c_1},
    (\Atom{\Obj A^1_1}{\Obj c_1} \land \Atom{\Obj A^1_0}{\Obj v_0}),
    \Atom{\Obj A^1_1}{\Obj c_1},\\
    \lforall[\Obj v_1][\Atom{\Obj A^1_0}{\Obj v_0}],
    \lexists[\Obj v_0][(\Atom{\Obj A^1_1}{\Obj c_1} \land \Atom{\Obj A^1_0}{\Obj v_0})]
}.
\end{multline*}
%
As can be seen from the second example, formation sequences
may contain ``junk'': \olup{formula}s which are redundant or do not
contribute to the construction.
\end{ex}

\begin{prop}\ollabel{prop:formed}
Every \olup{formula}~$A$ in~$\Frm[L]$ has a formation sequence.
\end{prop}

\begin{proof}
Suppose $A$ is atomic. Then the sequence $\tuple{A}$ is a
formation sequence for~$A$.
%
Now suppose that $B$ and~$C$ have formation sequences
$\tuple{B_0,\dotsc,B_n}$ and $\tuple{C_0,\dotsc,C_m}$
respectively.
%
\begin{enumerate}
  \tagitem{prvNot}{If $A \ident \lnot B$,
      then $\tuple{B_0,\dotsc,B_n,\lnot B_n}$
      is a formation sequence for~$A$.}{}
  \tagitem{prvAnd}{If $A \ident (B \land C)$,
      then $\tuple{B_0,\dotsc,B_n,C_0,\dotsc,C_m,(B_n \land C_m)}$
      is a formation sequence for~$A$.}{}
  \tagitem{prvOr}{If $A \ident (B \lor C)$,
      then $\tuple{B_0,\dotsc,B_n,C_0,\dotsc,C_m,(B_n \lor C_m)}$
      is a formation sequence for~$A$.}{}
  \tagitem{prvIf}{If $A \ident (B \lif C)$,
      then $\tuple{B_0,\dotsc,B_n,C_0,\dotsc,C_m,(B_n \lif C_m)}$
      is a formation sequence for~$A$.}{}
  \tagitem{prvIff}{If $A \ident (B \liff C)$,
      then $\tuple{B_0,\dotsc,B_n,C_0,\dotsc,C_m,(B_n \liff C_m)}$
      is a formation sequence for~$A$.}{}
  \tagitem{prvAll}{If $A \ident \lforall[x][B]$,
      then $\tuple{B_0,\dotsc,B_n,\lforall[x][B_n]}$
      is a formation sequence for~$A$.}{}
  \tagitem{prvEx}{If $A \ident \lexists[x][B]$,
      then $\tuple{B_0,\dotsc,B_n,\lexists[x][B_n]}$
      is a formation sequence for~$A$.}{}
\end{enumerate}
By the principle of induction on \olup{formula}s,
every \olup{formula} has a formation sequence.
\end{proof}

We can also prove the converse. This is important because it shows
that our two ways of defining formulas are equivalent: they give
the same results. It also means that we can prove theorems about
formulas by using ordinary induction on the length of formation
sequences.

\begin{lem}
\ollabel{lem:fseq-init}
Suppose that $\tuple{A_0,\dotsc,A_n}$ is a formation sequence
for~$A_n$, and that $k \leq n$. Then $\tuple{A_0,\dotsc,A_k}$
is a formation sequence for~$A_k$.
\end{lem}

\begin{proof}
Exercise.
\end{proof}

\begin{prob}
Prove \olref[fol][syn][fseq]{lem:fseq-init}.
\end{prob}

\begin{thm}
\ollabel{thm:fseq-frm-equiv}
$\Frm[L]$ is the set of all $\Lang L$-strings~$A$ such that
there exists a formula formation sequence for~$A$.
\end{thm}

\begin{proof}
Let $F$ be the set of all strings of symbols in the language~$\Lang L$
that have a formation sequence. We have seen in
\olref[fol][syn][fseq]{prop:formed} that $\Frm[L] \subseteq F$, so now
we prove the converse.

Suppose $A$ has a formation sequence $\tuple{A_0,\dotsc,A_n}$.
We prove that $A \in \Frm[L]$ by strong induction on~$n$.
Our induction hypothesis is that every string of symbols with a
formation sequence of length $m < n$ is in $\Frm[L]$.
By the definition of a formation sequence, either $A \ident A_n$ is
atomic or there must exist $j,k < n$ such that one of the
following is the case:
\begin{enumerate}
    \tagitem{prvNot}{$A \ident \lnot A_j$.}{}%
    \tagitem{prvAnd}{$A \ident (A_j \land A_k)$.}{}%
    \tagitem{prvOr}{$A \ident (A_j \lor A_k)$.}{}%
    \tagitem{prvIf}{$A \ident (A_j \lif A_k)$.}{}%
    \tagitem{prvIff}{$A \ident (A_j \liff A_k)$.}{}%
    \tagitem{prvAll}{$A \ident \lforall[x][A_j]$.}{}%
    \tagitem{prvEx}{$A \ident \lexists[x][A_j]$.}{}%
\end{enumerate}
Now we reason by cases. If $A$ is atomic then
$A_n \in \Frm[L_0]$. Suppose instead that
$A \equiv (A_j \land A_k)$. By 
\olref[fol][syn][fseq]{lem:fseq-init},
$\tuple{A_0,\dotsc,A_j}$ and $\tuple{A_0,\dotsc,A_k}$ are
formation sequences for $A_j$ and~$A_k$, respectively. Since
these are proper initial subsequences of the formation sequence
for~$A$, they both have length less than~$n$. Therefore by
the induction hypothesis, $A_j$ and~$A_k$ are in~$\Frm[L_0]$,
and by the definition of \olup{formula}, so is
$(A_j \land A_k)$. The other cases follow by parallel
reasoning.
\end{proof}

Formation sequences for terms have similar properties to those
for \olup{formula}s.

\begin{prop}
\ollabel{prop:fseq-trm-equiv}
$\Trm[L]$ is the set of all $\Lang L$-strings $t$
such that there exists a term formation sequence for~$t$.
\end{prop}

\begin{proof}
Exercise.
\end{proof}

\begin{prob}
Prove \olref[fol][syn][fseq]{prop:fseq-trm-equiv}.
Hint: use a similar strategy to that used in the proof of
\olref[fol][syn][fseq]{thm:fseq-frm-equiv}.
\end{prob}

There are two types of ``junk'' that can appear in formation
sequences: repeated elements, and elements that are irrelevant
to the construction of the formation or term. We can eliminate
both by looking at minimal formation sequences.

\begin{defn}[Minimal formation sequences]
\ollabel{defn:minimal-fseq}
A formation sequence $\tuple{A_0, \ldots, A_n}$ for a
formula~$A$ is a \emph{minimal formation sequence} for~$A$
if for every other formation sequence~$s$ for~$A$,
the length of~$s$ is greater than or equal to~$n+1$.

Similarly, a formation sequence $\tuple{t_0, \ldots, t_n}$
for a term~$t$ is a \emph{minimal formation sequence}
for~$t$ if for every other formation sequence~$s$ for~$t$,
the length of~$s$ is greater than or equal to~$n+1$.
\end{defn}

Note that a formula or term can have more than one minimal
formation sequence, but they will contain exactly the same
strings.

\begin{prop}
\ollabel{prop:subformula-equivs}
The following are equivalent:
\begin{enumerate}
    \item $B$ is a sub-\olup{formula} of~$A$.
    \item $B$ occurs in every formation sequence of~$A$.
    \item $B$ occurs in a minimal formation sequence of~$A$.
\end{enumerate}
\end{prop}

\begin{proof}
Exercise.
\end{proof}

\begin{prob}
Prove \olref[fol][syn][fseq]{prop:subformula-equivs}.
\end{prob}

\begin{history}
Formation sequences were introduced by Raymond Smullyan in his
textbook \emph{First-Order Logic} \citep{Smullyan1968}.
Additional properties of formation sequences were established by
\citet{Zuckerman1973}.
\end{history}


% ==== OLP file: content/first-order-logic/syntax-and-semantics/free-vars-sentences.tex ====
% Part: first-order-logic
% Chapter: syntax-and-semantics
% Section: free-vars-sentences



\olfileid{fol}{syn}{fvs}

\olsection{Free \printtoken{P}{variable} and \printtoken{P}{sentence}}

\begin{defn}[Free occurrences of \olup{variable}]
\ollabel{defn:free-occ}
The \emph{free} occurrences of \olup{variable} in \olup{formula} are defined
inductively as follows:
\begin{enumerate}
\item \indcase*{A}{$A$ is atomic}{all \olup{variable} occurrences in
  $\indfrm$ are free.}

\tagitem{prvNot}{\indcase{A}{\lnot B}{the free \olup{variable}
  occurrences of $\indfrm$ are exactly those of $B$.}}{}

\item \indcase{A}{(B \ast C)}{the free
  \olup{variable} occurrences of $\indfrm$ are those in $B$
  together with those in~$C$.}

\tagitem{prvAll}{\indcase{A}{\lforall[x][B]}{the free \olup{variable}
  occurrences in $\indfrm$ are all of those in~$B$ except for
  occurrences of~$x$.}}{}

\tagitem{prvEx}{\indcase{A}{\lexists[x][B]}{the free \olup{variable}
  occurrences in $\indfrm$ are all of those in~$B$ except for
  occurrences of~$x$.}}{}
\end{enumerate}
\end{defn}

\begin{defn}[Bound Variables]
An occurrence of \olup{variable} in a formula~$A$ is \emph{bound} if
it is not free.
\end{defn}

\begin{prob}
Give an inductive definition of the bound variable occurrences along
the lines of \olref[fol][syn][fvs]{defn:free-occ}.
\end{prob}

\begin{defn}[Scope]
\iftag{prvAll}{If $\lforall[x][B]$ is an occurrence of a subformula
  in a formula~$A$, then the corresponding occurrence of~$B$ in~$A$
  is called the \emph{scope} of the corresponding occurrence
  of~$\lforall[x]$. \iftag{prvEx}{Similarly for $\lexists[x]$.}{}}{If
  $\lexists[x][B]$ is an occurrence of a subformula in a
  formula~$A$, then the corresponding occurrence of~$B$ in~$A$ is
  called the \emph{scope} of the corresponding occurrence
  of~$\lexists[x]$.}

If $B$ is the scope of a quantifier occurrence
\iftag{prvAll}{$\lforall[x]$\iftag{prvEx}{ or
    $\lexists[x]$}{}}{$\lexists[x]$} in~$A$, then the free occurrences of
$x$ in~$B$ are bound in \iftag{prvAll}{$\lforall[x][B]$\iftag{prvEx}{ and
$\lexists[x][B]$}{}}{$\lexists[x][B]$}. We say that these
occurrences are \emph{bound by} the
mentioned quantifier occurrence.
\end{defn}

\begin{ex}
Consider the following formula:
\[
\lexists[\Obj v_0][\underbrace{\Atom{\Obj A^2_0}{\Obj v_0,\Obj v_1}}_{B}] 
\]
$B$ represents the scope of $\lexists[\Obj v_0]$. 
The quantifier binds the occurrence of $\Obj v_0$ in $B$, but
does not bind the occurrence of $\Obj v_1$. So $\Obj v_1$ is
a free variable in this case.


We can now see how this might work in a more complicated 
\olup{formula}~$A$:
\[
\lforall[\Obj v_0][\underbrace{(\Atom{\Obj A^1_0}{\Obj v_0} \lif
    \Atom{\Obj A^2_0}{\Obj v_0, \Obj v_1})}_{B}] \lif \lexists[\Obj
  v_1][\underbrace{(\Atom{\Obj A^2_1}{\Obj v_0, \Obj v_1} \lor \lforall[\Obj v_0][\overbrace{\lnot \Atom{\Obj A^1_1}{\Obj v_0}}^{D}])}_{C}]
\]
$B$ is the scope of the first $\lforall[\Obj v_0]$, $C$ is the scope
of $\lexists[\Obj v_1]$, and $D$ is the scope of the second
$\lforall[\Obj v_0]$.  The first $\lforall[\Obj v_0]$ binds the
occurrences of $\Obj v_0$ in~$B$, $\lexists[\Obj v_1]$ binds the occurrence
of $\Obj v_1$ in $C$, and the second $\lforall[\Obj v_0]$ binds the
occurrence of $\Obj v_0$ in~$D$.  The first occurrence of $\Obj v_1$
and the fourth occurrence of $\Obj v_0$ are free in~$A$. The last
occurrence of $\Obj v_0$ is free in $D$, but bound in $C$ and~$A$.
\end{ex}

\begin{defn}[Sentence]
!!^a{formula}~$A$ is \article{sentence} \emph{\olup{sentence}} iff it
contains no free occurrences of \olup{variable}s.
\end{defn}

% add examples!


% ==== OLP file: content/first-order-logic/syntax-and-semantics/intro-semantics.tex ====
% Part: first-order-logic
% Chapter: semantics
% Section: intro-semantics



\olfileid{fol}{syn}{its}

\olsection{Introduction}

Giving the meaning of expressions is the domain of semantics.  The
central concept in semantics is that of satisfaction in
\olup{structure}. !!^a{structure} gives meaning to the building blocks
of the language: \olup{domain} is a non-empty set of objects. The
quantifiers are interpreted as ranging over this domain, \olup{constant}s
are assigned elements in the domain, \olup{function}s are assigned
functions from the \olup{domain} to itself, and \olup{predicate}s are
assigned relations on the \olup{domain}.  The \olup{domain} together with
assignments to the basic vocabulary constitutes \olup{structure}.
!!^{variable}s may appear in \olup{formula}s, and in order to give a
semantics, we also have to assign \olup{element}s of the \olup{domain} to
them---this is a variable assignment. The satisfaction relation,
finally, brings these together. !!^a{formula} may be satisfied in
\olup{structure}~$\Struct{M}$ relative to \olup{variable} assignment~$s$,
written as $\Sat{M}{A}[s]$. This relation is also defined by
induction on the structure of~$A$, using the truth tables for the
logical connectives to define, say, satisfaction of $(A \land B)$ in
terms of satisfaction (or not) of $A$ and ~$B$.  It then turns out
that the \olup{variable} assignment is irrelevant if the \olup{formula}~$A$
is \olup{sentence}, i.e., has no free variables, and so we can talk of
\olup{sentence}s being simply satisfied (or not) in \olup{structure}s.

On the basis of the satisfaction relation $\Sat{M}{A}$ for \olup{sentence}s
we can then define the basic semantic notions of validity, entailment,
and satisfiability. !!^a{sentence} is valid, $\Entails A$, if every
\olup{structure} satisfies it. It is entailed by a set of \olup{sentence}s,
$\Gamma \Entails A$, if every \olup{structure} that satisfies all the
\olup{sentence}s in~$\Gamma$ also satisfies~$A$. And a set of \olup{sentence}s
is satisfiable if some \olup{structure} satisfies all \olup{sentence}s in it
at the same time.  Because \olup{formula}s are inductively defined, and
satisfaction is in turn defined by induction on the structure of
\olup{formula}s, we can use induction to prove properties of our
semantics and to relate the semantic notions defined.


% ==== OLP file: content/first-order-logic/syntax-and-semantics/intro-syntax.tex ====
% Part: first-order-logic
% Chapter: syntax
% Section: intro-syntax



\olfileid{fol}{syn}{itx}

\olsection{Introduction}

In order to develop the theory and metatheory of first-order logic, we
must first define the syntax and semantics of its expressions.  The
expressions of first-order logic are terms and \olup{formula}s.  Terms
are formed from \olup{variable}s, \olup{constant}s, and \olup{function}s.
!!^{formula}s, in turn, are formed from \olup{predicate}s together with
terms (these form the smallest, ``atomic'' \olup{formula}s), and then
from atomic \olup{formula}s we can form more complex ones using logical
connectives and quantifiers.  There are many different ways to set
down the formation rules; we give just one possible one. Other systems
will chose different symbols, will select different sets of
connectives as primitive, will use parentheses differently (or even not
at all, as in the case of so-called Polish notation).  What all
approaches have in common, though, is that the formation rules define
the set of terms and \olup{formula}s \emph{inductively}. If done
properly, every expression can result essentially in only one way
according to the formation rules.  The inductive definition resulting
in expressions that are \emph{uniquely readable} means we can give
meanings to these expressions using the same method---inductive
definition.


% ==== OLP file: content/first-order-logic/syntax-and-semantics/introduction.tex ====
% Part: first-order-logic
% Chapter: syntax-and-semantics
% Section: introduction



\olfileid{fol}{syn}{int}

\olsection{Introduction}

In order to develop the theory and metatheory of first-order logic, we
must first define the syntax and semantics of its expressions.  The
expressions of first-order logic are terms and \olup{formula}s.  Terms
are formed from \olup{variable}s, \olup{constant}s, and \olup{function}s.
!!^{formula}s, in turn, are formed from \olup{predicate}s together with
terms (these form the smallest, ``atomic'' \olup{formula}s), and then
from atomic \olup{formula}s we can form more complex ones using logical
connectives and quantifiers.  There are many different ways to set
down the formation rules; we give just one possible one. Other systems
will chose different symbols, will select different sets of
connectives as primitive, will use parentheses differently (or even not
at all, as in the case of so-called Polish notation).  What all
approaches have in common, though, is that the formation rules define
the set of terms and \olup{formula}s \emph{inductively}. If done
properly, every expression can result essentially in only one way
according to the formation rules.  The inductive definition resulting
in expressions that are \emph{uniquely readable} means we can give
meanings to these expressions using the same method---inductive
definition.

Giving the meaning of expressions is the domain of semantics.  The
central concept in semantics is that of satisfaction in
\olup{structure}. !!^a{structure} gives meaning to the building blocks
of the language: \olup{domain} is a non-empty set of objects. The
quantifiers are interpreted as ranging over this domain, \olup{constant}s
are assigned elements in the domain, \olup{function}s are assigned
functions from the \olup{domain} to itself, and \olup{predicate}s are
assigned relations on the \olup{domain}.  The \olup{domain} together with
assignments to the basic vocabulary constitutes
\olup{structure}. !!^{variable}s may appear in \olup{formula}s, and in
order to give a semantics, we also have to assign \olup{element}s of the
\olup{domain} to them---this is a variable assignment. The satisfaction
relation, finally, brings these together. !!^a{formula} may be satisfied
in \olup{structure}~$\Struct{M}$ relative to a variable assignment~$s$,
written as $\Sat{M}{A}[s]$. This relation is also defined by
induction on the structure of~$A$, using the truth tables for the
logical connectives to define, say, satisfaction of $A \land B$ in
terms of satisfaction (or not) of $A$ and ~$B$.  It then turns out
that the variable assignment is irrelevant if the \olup{formula}~$A$ is
\olup{sentence}, i.e., has no free variables, and so we can talk of
\olup{sentence}s being simply satisfied (or not) in \olup{structure}s.

On the basis of the satisfaction relation $\Sat{M}{A}$ for sentences
we can then define the basic semantic notions of validity, entailment,
and satisfiability.  A sentence is valid, $\Entails A$, if every
structure satisfies it. It is entailed by a set of \olup{sentence}s,
$\Gamma \Entails A$, if every \olup{structure} that satisfies all the
\olup{sentence}s in~$\Gamma$ also satisfies~$A$. And a set of sentences
is satisfiable if some \olup{structure} satisfies all \olup{sentence}s in it
at the same time.  Because \olup{formula}s are inductively defined, and
satisfaction is in turn defined by induction on the structure of
\olup{formula}s, we can use induction to prove properties of our
semantics and to relate the semantic notions defined.





% ==== OLP file: content/first-order-logic/syntax-and-semantics/main-operator.tex ====
% Part: first-order-logic
% Chapter: syntax-and-semantics
% Section: main-operator



\olfileid{fol}{syn}{mai}

\olsection{\printtoken{S}{main operator} of a Formula}

\begin{explain}
It is often useful to talk about the last operator used in
constructing \olup{formula}~$A$.  This operator is called the \emph{main
  operator} of~$A$. Intuitively, it is the ``outermost'' operator
of $A$. For example, the main operator of $\lnot A$ is $\lnot$,
the main operator of $(A \lor B)$ is $\lor$, etc.
\end{explain}


\begin{defn}[!!^{main operator}]
\ollabel{def:main-op}
The \emph{\olup{main operator}} of \olup{formula}~$A$ is
defined as follows:
\begin{enumerate}
\item \indcase*{A}{A}{$\indfrm$ has no \olup{main operator}.}

\tagitem{prvNot}{\indcase{A}{\lnot B}{the \olup{main operator} of $\indfrm$
  is~$\lnot$.}}{}

\tagitem{prvAnd}{\indcase{A}{(B \land C)}{the \olup{main operator} of
  $\indfrm$ is~$\land$.}}{}

\tagitem{prvOr}{\indcase{A}{(B \lor C)}{the \olup{main operator} of
  $\indfrm$ is~$\lor$.}}{}

\tagitem{prvIf}{\indcase{A}{(B \lif C)}{the \olup{main operator} of
  $\indfrm$ is~$\lif$.}}{}

\tagitem{prvIff}{\indcase{A}{(B \liff C)}{the \olup{main operator} of
  $\indfrm$ is~$\liff$.}}{}

\tagitem{prvAll}{\indcase{A}{\lforall[x][B]}{the \olup{main operator}
  of $\indfrm$ is~$\lforall$.}}{}

\tagitem{prvEx}{\indcase{A}{\lexists[x][B]}{the \olup{main operator} of
  $\indfrm$ is~$\lexists$.}}{}
\end{enumerate}
\end{defn}

In each case, we intend the specific indicated \emph{occurrence} of
the \olup{main operator} in the formula. For instance, since the formula
$((D \lif E) \lif (E \lif D))$ is of the form $(B \lif C)$ where
$B$ is $(D \lif E)$ and $C$ is $(E \lif D)$, the second
occurrence of $\lif$ is the \olup{main operator}.

\begin{explain}
This is a \emph{recursive} definition of a function which maps all
non-atomic \olup{formula}s to their \olup{main operator} occurrence. Because
of the way \olup{formula}s are defined inductively, every
\olup{formula}~$A$ satisfies one of the cases in \olref{def:main-op}.
This guarantees that for each non-atomic \olup{formula}~$A$ \olup{main
  operator} exists. Because each \olup{formula} satisfies only one of these
conditions, and because the smaller \olup{formula}s from which $A$ is
constructed are uniquely determined in each case, the \olup{main
  operator} occurrence of~$A$ is unique, and so we have defined a
function.
\end{explain}

We call \olup{formula}s by the names in \olref{tab:main-op} depending on
which symbol their \olup{main operator}
is.\iftag{defNot,defOr,defAnd,defIf,defIff,defTrue,defFalse,defEx,defAll}
{ Recall, however, that defined operators do not officially appear in
\olup{formula}s. They are just abbreviations, so officially they cannot
be the main operator of a formula. In proofs about all \olup{formula}s
they therefore do not have to be treated separately.}

\begin{table}[h]
\centering
\begin{tabular}{c | c | c}
!!^{main operator} & Type of \olup{formula} & Example\\
\hline
none & atomic (\olup{formula}) &
\iftag{prvFalse}{$\lfalse$,}{}
\iftag{prvTrue}{$\ltrue$,}{}
$\Atom{R}{t_1, \dots, t_n}$,
$\eq[t_1][t_2]$\\
$\lnot$ & negation & $\lnot A$ \\
$\land$ & conjunction & $(A \land B$) \\
$\lor$ & disjunction & $(A \lor B$) \\
$\lif$ & \olup{conditional} & $(A \lif B$) \\
$\liff$ & \olup{biconditional} & $(A \liff B)$ \\
$\lforall[][]$ & universal (\olup{formula})& $\lforall[x][A]$ \\
$\lexists[][]$ & existential (\olup{formula})& $\lexists[x][A]$
\end{tabular}
\caption{Main operator and names of \olup{formula}s}
\ollabel{tab:main-op}
\end{table}


% ==== OLP file: content/first-order-logic/syntax-and-semantics/satisfaction.tex ====
% Part: first-order-logic
% Chapter: syntax-and-semantics
% Section: satisfaction



\olfileid{fol}{syn}{sat}

\olsection{Satisfaction of \article{formula} \printtoken{S}{formula}
  in \article{structure} \printtoken{S}{structure}}

\begin{explain}
The basic notion that relates expressions such as terms and
\olup{formula}s, on the one hand, and \olup{structure}s on the other, are
those of \emph{\olup{value}} of a term and \emph{satisfaction} of 
\olup{formula}.  Informally, the \olup{value} of a term is an \olup{element} of
\olup{structure}---if the term is just a constant, its \olup{value} is the
object assigned to the constant by the \olup{structure}, and if it is
built up using \olup{function}s, the \olup{value} is computed from the
\olup{value}s of constants and the functions assigned to the functions in
the term.  !!^a{formula} is \emph{satisfied} in \olup{structure} if the
interpretation given to the predicates makes the \olup{formula} true in
the domain of the \olup{structure}. This notion of satisfaction is
specified inductively: the specification of the \olup{structure} directly
states when atomic \olup{formula}s are satisfied, and we define when a
complex \olup{formula} is satisfied depending on the main connective or
quantifier and whether or not the immediate \olup{subformula}s are
satisfied. 

The case of the quantifiers here is a bit tricky, as the
immediate \olup{subformula} of a quantified \olup{formula} has a free
\olup{variable}, and \olup{structure}s don't specify the \olup{value}s of
\olup{variable}s.  In order to deal with this difficulty, we also
introduce \emph{variable assignments} and define satisfaction not with
respect to \olup{structure} alone, but with respect to \olup{structure}
plus \olup{variable} assignment.
\end{explain}

\begin{defn}[Variable Assignment]
A \emph{variable assignment}~$s$ for \olup{structure}~$\Struct{M}$ is a
function which maps each \olup{variable} to \olup{element} of~$\Domain M$,
i.e., $s\colon \Var \to \Domain M$.
\end{defn}

\begin{explain}
!!^a{structure} assigns \olup{value} to each \olup{constant}, and a
variable assignment to each variable.  But we want to use terms built
up from them to also name \olup{element}s of the \olup{domain}.  For this we
define the \olup{value} of terms inductively. For \olup{constant}s and
variables the value is just as the \olup{structure} or the variable
assignment specifies it; for more complex terms it is computed
recursively using the functions the \olup{structure} assigns to the
\olup{function}s.
\end{explain}

\begin{defn}[!!^{value} of Terms]
If $t$ is a term of the language~$\Lang L$, $\Struct M$ is a
\olup{structure} for~$\Lang L$, and $s$ is \olup{variable} assignment
for~$\Struct M$, the \emph{\olup{value}}~$\Value{t}{M}[s]$ is defined as
follows:
\begin{enumerate}
\item \indcase{t}{c}{$\Value{\indfrm}{M}[s] = \Assign{\indcomplex}{M}$.}
\item \indcase{t}{x}{$\Value{\indfrm}{M}[s] = s(\indcomplex)$.}
\item \indcase{t}{\Atom{f}{t_1, \ldots, t_n}}{
\[
\Value{\indfrm}{M}[s] = \Assign{f}{M}(\Value{t_1}{M}[s], \ldots,
\Value{t_n}{M}[s]).
\]}
\end{enumerate}
\end{defn}

\begin{defn}[$x$-Variant]
If $s$ is \olup{variable} assignment for \olup{structure}~$\Struct M$, then any
\olup{variable} assignment~$s'$ for~$\Struct M$ which differs from~$s$ at most
in what it assigns to~$x$ is called an \emph{$x$-variant} of~$s$.  If
$s'$ is an $x$-variant of~$s$ we write $\varAssign{s'}{s}{x}$.
\end{defn}

\begin{explain}
Note that an $x$-variant of an assignment~$s$ does not \emph{have} to
assign something different to~$x$.  In fact, every assignment counts
as an $x$-variant of itself.
\end{explain}

\begin{defn}
  If $s$ is \olup{variable} assignment for \olup{structure}~$\Struct M$
  and $m \in \Domain{M}$, then the assignment~$\Subst{s}{m}{x}$ is the
  variable assignment defined by
  \[\Subst{s}{m}{x}(y) = \begin{cases}
    m & \text{if } y \ident x\\
    s(y) & \text{otherwise}.
  \end{cases}\]
\end{defn}

In other words, $\Subst{s}{m}{x}$ is the particular $x$-variant of~$s$
which assigns the domain \olup{element}~$m$ to~$x$, and assigns the same
things to \olup{variable}s other than~$x$ that $s$ does.

\begin{defn}[Satisfaction]
\ollabel{defn:satisfaction}
Satisfaction of \olup{formula}~$A$ in \olup{structure}~$\Struct M$
relative to \olup{variable} assignment~$s$, in symbols:
$\Sat{M}{A}[s]$, is defined recursively as follows. (We write
$\Sat/{M}{A}[s]$ to mean ``not $\Sat{M}{A}[s]$.'')
\begin{enumerate}
\tagitem{prvFalse}{%
  \indcase{A}{\lfalse}{$\Sat/{M}{\indfrm}[s]$.}}{}

\tagitem{prvTrue}{%
  \indcase{A}{\ltrue}{$\Sat{M}{\indfrm}[s]$.}}{}

\item \indcase{A}{\Atom{R}{t_1, \dots, t_n}}{$\Sat{M}{\indfrm}[s]$
  iff $\langle \Value{t_1}{M}[s], \dots, \Value{t_n}{M}[s] \rangle \in
  \Assign{R}{M}$.}

\item \indcase{A}{\eq[t_1][t_2]}{$\Sat{M}{\indfrm}[s]$ iff
  $\Value{t_1}{M}[s] = \Value{t_2}{M}[s]$.}

\tagitem{prvNot}{%
  \indcase{A}{\lnot B}{$\Sat{M}{\indfrm}[s]$ iff
    $\Sat/{M}{B}[s]$.}}{}

\tagitem{prvAnd}{%
  \indcase{A}{(B \land C)}{$\Sat{M}{\indfrm}[s]$ iff $\Sat{M}{B}[s]$
    and $\Sat{M}{C}[s]$.}}{}

\tagitem{prvOr}{%
  \indcase{A}{(B \lor C)}{$\Sat{M}{\indfrm}[s]$ iff
    $\Sat{M}{B}[s]$ or $\Sat{M}{C}[s]$ (or both).}}{}

\tagitem{prvIf}{%
  \indcase{A}{(B \lif C)}{$\Sat{M}{\indfrm}[s]$ iff $\Sat/{M}{B}[s]$
    or $\Sat{M}{C}[s]$ (or both).}}{}

\tagitem{prvIff}{%
  \indcase{A}{(B \liff C)}{$\Sat{M}{\indfrm}[s]$ iff either both
    $\Sat{M}{B}[s]$ and $\Sat{M}{C}[s]$, or neither $\Sat{M}{B}[s]$
    nor $\Sat{M}{C}[s]$.}}{}

\tagitem{prvAll}{%
  \indcase{A}{\lforall[x][B]}{$\Sat{M}{\indfrm}[s]$ iff for every
    \olup{element}~$m \in \Domain M$, $\Sat{M}{B}[\Subst{s}{m}{x}]$.}}{}

\tagitem{prvEx}{%
  \indcase{A}{\lexists[x][B]}{$\Sat{M}{\indfrm}[s]$ iff for at least
  one \olup{element}~$m \in \Domain M$, $\Sat{M}{B}[\Subst{s}{m}{x}]$.}}{}
\end{enumerate}
\end{defn}

\begin{explain}
The variable assignments are important in the last
\iftag{notprvEx,notprvAll}{clause}{two clauses}.\iftag{prvAll}{ We
cannot define satisfaction of $\lforall[x][B(x)]$ by ``for all $m \in
\Domain{M}$, $\Sat{M}{B(m)}$.''}{}\iftag{prvEx}{ We cannot define
satisfaction of $\lexists[x][B(x)]$ by ``for at least one $m \in
\Domain{M}$, $\Sat{M}{B(m)}$.''}{} The reason is that if $m \in
\Domain M$, it is not a symbol of the language, and so $B(m)$~is not
\olup{formula} (that is, $\Subst{B}{m}{x}$ is undefined).  We also
cannot assume that we have \olup{constant}s or terms available that name
every \olup{element} of~$\Struct{M}$, since there is nothing in the
definition of \olup{structure}s that requires it.  In the standard
language, the set of \olup{constant}s is \olup{denumerable}, so if
$\Domain{M}$ is not \olup{enumerable} there aren't even enough
\olup{constant}s to name every object. 

We solve this problem by introducing \olup{variable} assignments, which
allow us to link variables directly with \olup{element}s of the domain.
Then instead of saying that, e.g.,
\iftag{prvEx}{$\lexists[x][B(x)]$}{$\lforall[x][B(x)]$} is satisfied
in~$\Struct M$ iff \iftag{prvEx}{for at least one $m \in
\Domain{M}$}{for all $m \in \Domain{M}$, $\Sat{M}{B(m)}$}, we say it
is satisfied in~$\Struct M$ \emph{relative to}~$s$ iff $B(x)$ is
satisfied relative to~$\Subst{s}{m}{x}$ \iftag{prvEx}{for at least
one}{for every} $m \in \Domain M$.
\end{explain}

\begin{ex}
Let $\Lang{L} = \{a, b, f, R\}$ where $a$ and $b$ are \olup{constant}s,
$f$~is a two-place \olup{function}, and $R$~is a two-place \olup{predicate}.
Consider the \olup{structure}~$\Struct{M}$ defined by:
\begin{enumerate}
\item $\Domain M = \{1, 2, 3, 4\}$
\item $\Assign{a}{M} = 1$
\item $\Assign{b}{M} = 2$
\item $\Assign{f}{M}(x, y) = x+y$ if $x+y \le 3$ and $= 3$ otherwise.
\item $\Assign{R}{M} = \{\tuple{1, 1}, \tuple{1, 2}, \tuple{2, 3}, \tuple{2, 4}\}$
\end{enumerate}
The function $s(x) = 1$ that assigns $1 \in \Domain{M}$ to every
\olup{variable} is a variable assignment for~$\Struct{M}$.

Then
\begin{align*}
\Value{f(a,b)}{M}[s] & = \Assign{f}{M}(\Value{a}{M}[s], \Value{b}{M}[s]).
\intertext{Since $a$ and $b$ are \olup{constant}s, $\Value{a}{M}[s]
  = \Assign{a}{M} = 1$ and $\Value{b}{M}[s] = \Assign{b}{M} = 2$. So}
\Value{f(a,b)}{M}[s] & = \Assign{f}{M}(1, 2) = 1+2 = 3.
\intertext{To compute the value of $f(f(a,b),a)$ we have to consider}
\Value{f(f(a,b),a)}{M}[s] & = \Assign{f}{M}(\Value{f(a, b)}{M}[s],
\Value{a}{M}[s]) = \Assign{f}{M}(3, 1) = 3,
\intertext{since $3+1 > 3$. Since $s(x) = 1$ and $\Value{x}{M}[s] =
  s(x)$, we also have}
\Value{f(f(a,b),x)}{M}[s] & = \Assign{f}{M}(\Value{f(a, b)}{M}[s],
\Value{x}{M}[s]) = \Assign{f}{M}(3, 1) = 3,
\end{align*}

An atomic \olup{formula}~$R(t_1, t_2)$ is satisfied if the tuple of
values of its arguments, i.e., $\tuple{\Value{t_1}{M}[s],
  \Value{t_2}{M}[s]}$, is \olup{element} of~$\Assign{R}{M}$. So, e.g., we
have $\Sat{M}{R(b,f(a,b))}[s]$ since $\tuple{\Value{b}{M},
  \Value{f(a,b)}{M}} = \tuple{2, 3} \in \Assign{R}{M}$, but
$\Sat/{M}{R(x, f(a,b))}[s]$ since $\tuple{1, 3} \notin \Assign{R}{M}[s]$.

To determine if a non-atomic formula~$A$ is satisfied, you apply the
clauses in the inductive definition that applies to the main
connective. For instance, the main connective in $R(a, a) \lif (R(b,
x) \lor R(x, b))$ is the~$\lif$, and
\begin{align*}
  & \Sat{M}{R(a, a) \lif (R(b, x) \lor R(x, b))}[s] \text{ iff }\\
  & \qquad
  \Sat/{M}{R(a,a)}[s] \text{ or } \Sat{M}{R(b, x) \lor R(x, b)}[s]
  \intertext{Since $\Sat{M}{R(a,a)}[s]$ (because $\tuple{1,1} \in
    \Assign{R}{M}$) we can't yet determine the answer and must first
    figure out if $\Sat{M}{R(b, x) \lor R(x, b)}[s]$:}
  & \Sat{M}{R(b, x) \lor R(x, b)}[s] \text{ iff }\\
  & \qquad \Sat{M}{R(b,x)}[s] \text{ or } \Sat{M}{R(x, b)}[s]
  \intertext{And this is the case, since $\Sat{M}{R(x, b)}[s]$
    (because $\tuple{1,2} \in \Assign{R}{M}$).}
\end{align*}

Recall that an $x$-variant of~$s$ is a variable assignment that
differs from $s$ at most in what it assigns to~$x$. For every
\olup{element} of~$\Domain{M}$, there is an $x$-variant of~$s$:
\begin{align*}
  s_1 & = \Subst{s}{1}{x}, &
  s_2 & = \Subst{s}{2}{x},\\
  s_3 & = \Subst{s}{3}{x}, &
  s_4 & = \Subst{s}{4}{x}.
\end{align*}
So, e.g., $s_2(x) = 2$ and $s_2(y) = s(y) = 1$ for all variables~$y$
other than~$x$. These are all the $x$-variants of~$s$ for the
structure~$\Struct{M}$, since $\Domain{M} = \{1, 2, 3, 4\}$. Note, in
particular, that $s_1 = s$ ($s$~is always an $x$-variant of itself).

\iftag{prvEx}{To determine if an existentially quantified
  \olup{formula}~$\lexists[x][A(x)]$ is satisfied, we have to determine
  if $\Sat{M}{A(x)}[\Subst{s}{m}{x}]$ for at least one $m \in \Domain
  M$. So,
  \[
  \Sat{M}{\lexists[x][(R(b,x) \lor R(x,b))]}[s],
  \]
  since $\Sat{M}{R(b,x) \lor R(x, b)}[\Subst{s}{1}{x}]$
  ($\Subst{s}{3}{x}$ would also fit the bill).  But,
  \[
  \Sat/{M}{\lexists[x][(R(b,x) \land R(x,b))]}[s]
  \]
  since, whichever $m \in \Domain{M}$ we pick, $\Sat/{M}{R(b,x) \land R(x,b)}[\Subst{s}{m}{x}]$.}{}

\iftag{prvAll}{To determine if a universally quantified
  \olup{formula}~$\lforall[x][A(x)]$ is satisfied, we have to determine
  if $\Sat{M}{A(x)}[\Subst{s}{m}{x}]$ for all $m \in \Domain M$. So,
  \[
  \Sat{M}{\lforall[x][(R(x,a) \lif R(a,x))]}[s],
  \]
  since $\Sat{M}{R(x,a) \lif R(a,x)}[\Subst{s}{m}{x}]$ for all $m \in
  \Domain M$. For $m = 1$, we have $\Sat{M}{R(a,x)}[\Subst{s}{1}{x}]$
  so the consequent is true; for $m = 2$, $3$, and~$4$, we have
  $\Sat/{M}{R(x,a)}[\Subst{s}{m}{x}]$, so the antecedent is false.
  But,
  \[
  \Sat/{M}{\lforall[x][(R(a,x) \lif R(x,a))]}[s]
  \]
  since $\Sat/{M}{R(a,x) \lif R(x,a)}[\Subst{s}{2}{x}]$ (because
  $\Sat{M}{R(a, x)}[\Subst{s}{2}{x}]$ and $\Sat/{M}{R(x,
  a)}[\Subst{s}{2}{x}]$).}{}

\iftag{defEx}{To determine if an existentially quantified
  \olup{formula}~$\lexists[x][A(x)]$ is satisfied, we have to determine
  if $\Sat{M}{\lnot \lforall[x][\lnot A(x)]}[s]$. For instance, we
  have
  \[
  \Sat{M}{\lexists[x][(R(b,x) \lor R(x,b))]}[s].
  \]
  First, $\Sat{M}{R(b,x) \lor R(x, b)}[\Subst{s}{1}{x}]$
  ($\Subst{s}{3}{x}$ would also fit the bill).  So,
  $\Sat/{M}{\lnot(R(b,x) \lor R(x,b))}[\Subst{s}{1}{x}]$, thus
  $\Sat/{M}{\lforall[x][\lnot((R(b,x) \lor R(x,b))]}[s]$, and
  therefore $\Sat{M}{\lnot\lforall[x][\lnot((R(b,x) \lor
  R(x,b))]}[s]$. On the other hand,
  \[
  \Sat/{M}{\lexists[x][(R(b,x) \land R(x,b))],}[s].
  \]
  That's because $\Sat{M}{\lforall[x][\lnot(R(b,x) \land
  R(x,b))]}[s]$, since for no $m \in \Domain M$, $\Sat{M}{R(b,x) \land
  R(x,b)}[\Subst{s}{m}{x}]$. As you can probably guess from these
  examples, $\Sat{M}{\lexists[x][A(x)]}[s]$ iff
  $\Sat{M}{A(x)}[\Subst{s}{m}{x}]$ for at least one $m \in \Domain
  M$.}{}

\iftag{defAll}{To determine if a universally quantified
  \olup{formula}~$\lforall[x][A(x)]$ is satisfied, we have to determine
  if $\Sat{M}{\lnot\lexists[x][\lnot A(x)]}[s]$. For instance,
  \[
  \Sat{M}{\lforall[x][(R(x,a) \lif R(a,x))]}[s],
  \]
  First, $\Sat{M}{R(x,a) \lif R(a,x)}[\Subst{s}{m}{x}]$ for all $m \in
  \Domain M$ ($\Sat{M}{R(a,x)}[\Subst{s}{1}{x}]$ and
  $\Sat/{M}{R(a,x)}[\Subst{s}{m}{x}]$ for $m = 2$, $3$, or~$4$).
  Thus, there is no $m \in \Domain M$ such that
  $\Sat{M}{\lnot(R(x,a) \lif R(a,x))}[\Subst{s}{m}{x}]$ and hence
  $\Sat/{M}{\lexists[x][\lnot(R(x,a) \lif R(a,x))]}[s]$. Therefore,
  $\Sat{M}{\lnot\lexists[x][\lnot(R(x,a) \lif R(a,x))]}[s]$. On the
  other hand,
  \[
  \Sat/{M}{\lforall[x][(R(a,x) \lif R(x,a))]}[s],
  \]
  since $\Sat/{M}{R(a,x) \lif R(x,a)}[\Subst{s}{2}{x}]$, and so
  $\Sat{M}{\lexists[x][\lnot(R(a,x) \lif R(x,a))]}[s]$. As you can
  probably guess from these examples, $\Sat{M}{\lforall[x][A(x)]}[s]$
  iff $\Sat{M}{A(x)}[\Subst{s}{m}{x}]$ for every $m \in \Domain
  M$.}{}

For a more complicated case, consider
\[
\lforall[x][(R(a,x) \lif \lexists[y][R(x,y)])].
\]
Since $\Sat/{M}{R(a,x)}[\Subst{s}{3}{x}]$ and
$\Sat/{M}{R(a,x)}[\Subst{s}{4}{x}]$, the interesting cases where we
have to worry about the consequent of the conditional are only $m = 1$
and $ = 2$. Does $\Sat{M}{\lexists[y][R(x,y)]}[\Subst{s}{1}{x}]$
hold? It does if there is at least one $n \in \Domain M$ so that
$\Sat{M}{R(x,y)}[\Subst{\Subst{s}{1}{x}}{n}{y}]$. In fact, if we take
$n = 1$, we have $\Subst{\Subst{s}{1}{x}}{n}{y} = \Subst{s}{1}{y} =
s$. Since $s(x) = 1$, $s(y) = 1$, and $\tuple{1,1} \in \Assign{R}{M}$,
the answer is yes. 

To determine if $\Sat{M}{\lexists[y][R(x,y)]}[\Subst{s}{2}{x}]$, we
have to look at the \olup{variable} assignments
$\Subst{\Subst{s}{2}{x}}{n}{y}$. Here, for $n = 1$, this assignment
is~$s_2 = \Subst{s}{2}{x}$, which does not satisfy $R(x,y)$ ($s_2(x) =
2$, $s_2(y) = 1$, and $\tuple{2,1}\notin \Assign{R}{M}$). However,
consider $\Subst{\Subst{s}{2}{x}}{3}{y} = \Subst{s_2}{3}{y}$.
$\Sat{M}{R(x,y)}[\Subst{s_2}{3}{y}]$ since $\tuple{2,3} \in
\Assign{R}{M}$, and so $\Sat{M}{\lexists[y][R(x,y)]}[s_2]$.

So, for all $n \in \Domain M$, either
$\Sat/{M}{R(a,x)}[\Subst{s}{m}{x}]$ (if $m = 3$, $4$) or
$\Sat{M}{\lexists[y][R(x,y)]}[\Subst{s}{m}{x}]$ (if $m = 1$, $2$), and so
\[
\Sat{M}{\lforall[x][(R(a,x) \lif \lexists[y][R(x,y)])]}[s].
\]
On the other hand,
\[
\Sat/{M}{\lexists[x][(R(a,x) \land \lforall[y][R(x,y)])]}[s].
\]
We have $\Sat{M}{R(a,x)}[\Subst{s}{m}{x}]$ only for $m = 1$ and $m =
2$. But for both of these values of~$m$, there is in turn an $n \in
\Domain M$, namely $n = 4$, so that
$\Sat/{M}{R(x,y)}[\Subst{\Subst{s}{m}{x}}{n}{y}]$ and so
$\Sat/{M}{\lforall[y][R(x,y)]}[\Subst{s}{m}{x}]$ for $m = 1$ and $m =
2$. In sum, there is no $m \in \Domain M$ such that $\Sat{M}{R(a,x)
\land \lforall[y][R(x,y)]}[\Subst{s}{m}{x}]$.
\end{ex}

\iftag{defEx}{%
\begin{prop}\ollabel{prop:sat-ex}
$\Sat{M}{\lexists[x][B(x)]}[s]$ iff there is an $x$-variant $s'$ of $s$
    so that $\Sat{M}{B(x)}[s']$.
\end{prop}

\begin{proof}
  Exercise.
\end{proof}
}{}

\tagprob{defEx}
\begin{prob}
  Prove \olref[fol][syn][sat]{prop:sat-ex}
\end{prob}
\tagendprob

\iftag{defAll}{%
\begin{prop}\ollabel{prop:sat-all}
$\Sat{M}{\lforall[x][B(x)]}[s]$ iff for every $x$-variant~$s'$ of $s$,
  $\Sat{M}{B(x)}[s']$
\end{prop}

\begin{proof}
  Exercise.
\end{proof}
}{}

\tagprob{defAll}
\begin{prob}
  Prove \olref[fol][syn][sat]{prop:sat-all}
\end{prob}
\tagendprob

\begin{prob}
Let $\Lang L = \{c, f, A\}$ with one \olup{constant}, one one-place
\olup{function} and one two-place \olup{predicate}, and let the
\olup{structure}~$\Struct{M}$ be given by
\begin{enumerate}
\item $\Domain M = \{1, 2, 3\}$
\item $\Assign{c}{M} = 3$
\item $\Assign{f}{M}(1) = 2, \Assign{f}{M}(2) = 3, \Assign{f}{M}(3) = 2$
\item $\Assign{A}{M} = \{\tuple{1, 2}, \tuple{2, 3}, \tuple{3, 3}\}$
\end{enumerate}
(a) Let $s(v) = 1$ for all \olup{variable}s~$v$.  Find out whether
\[
\Sat{M}{\lexists[x][(A(f(z), c) \lif \lforall[y][(A(y, x) \lor A(f(y),
      x))])]}[s]
\]
Explain why or why not.

(b) Give a different structure and \olup{variable} assignment in which the
\olup{formula} is not satisfied.
\end{prob}


% ==== OLP file: content/first-order-logic/syntax-and-semantics/semantic-notions.tex ====
% Part: first-order-logic
% Chapter: syntax-and-semantics
% Section: Semantic Notions



\olfileid{fol}{syn}{sem}
\olsection{Semantic Notions}

\begin{explain}
Given the definition of \olup{structure}s for first-order languages, we can
define some basic semantic properties of and relationships between
sentences.  The simplest of these is the notion of \emph{validity} of
a sentence.  A sentence is valid if it is satisfied in every
\olup{structure}.  Valid sentences are those that are satisfied regardless of
how the non-logical symbols in it are interpreted.  Valid sentences
are therefore also called \emph{logical truths}---they are true, i.e.,
satisfied, in any \olup{structure} and hence their truth depends only on the
logical symbols occurring in them and their syntactic \olup{structure}, but not
on the non-logical symbols or their interpretation.
\end{explain}

\begin{defn}[Validity]
A sentence $A$ is \emph{valid}, $\Entails A$, iff $\Sat{M}{A}$ for every
\olup{structure}~$\Struct M$.
\end{defn}

\begin{defn}[Entailment]
A set of sentences~$\Gamma$ \emph{entails} a sentence~$A$, $\Gamma
\Entails A$, iff for every \olup{structure}~$\Struct M$ with
$\Sat{M}{\Gamma}$, $\Sat{M}{A}$.
\end{defn}


\begin{defn}[Satisfiability]
A set of sentences~$\Gamma$ is \emph{satisfiable} if $\Sat{M}{\Gamma}$
for some \olup{structure}~$\Struct M$.  If $\Gamma$ is not satisfiable it is
called \emph{unsatisfiable}.
\end{defn}

\begin{prop}
A sentence $A$ is valid iff $\Gamma \Entails A$ for every set of
sentences~$\Gamma$.
\end{prop}

\begin{proof}
For the forward direction, let $A$ be valid, and let $\Gamma$ be a
set of sentences. Let $\Struct M$ be \olup{structure} so that
$\Sat{M}{\Gamma}$. Since $A$ is valid, $\Sat{M}{A}$, hence $\Gamma
\Entails A$.

For the contrapositive of the reverse direction, let $A$ be invalid,
so there is \olup{structure}~$\Struct M$ with $\Sat/{M}{A}$. When $\Gamma
= \{ \ltrue \}$, since $\ltrue$ is valid, $\Sat{M}{\Gamma}$. Hence,
there is \olup{structure}~$\Struct M$ so that $\Sat{M}{\Gamma}$ but
$\Sat/{M}{A}$, hence $\Gamma$ does not entail $A$.
\end{proof}

\begin{prop}
\ollabel{prop:entails-unsat}
$\Gamma \Entails A$ iff $\Gamma \cup \{\lnot A\}$ is unsatisfiable.
\end{prop}

\begin{proof}
For the forward direction, suppose $\Gamma \Entails A$ and suppose to the
contrary that there is \olup{structure}~$\Struct M$ so that $\Sat{M}{\Gamma
  \cup \{ \lnot A \}}$. Since $\Sat{M}{\Gamma}$ and $\Gamma \Entails
A$, $\Sat{M}{A}$. Also, since $\Sat{M}{\Gamma\cup \{ \lnot A \}}$,
$\Sat{M}{\lnot A}$, so we have both $\Sat{M}{A}$ and $\Sat/{M}{A}$,
a contradiction. Hence, there can be no such \olup{structure}~$\Struct M$, so
$\Gamma \cup \{ \lnot A \}$ is unsatisfiable.

For the reverse direction, suppose $\Gamma \cup \{ \lnot A \}$ is
unsatisfiable. So for every \olup{structure}~$\Struct M$, either
$\Sat/{M}{\Gamma}$ or $\Sat{M}{A}$. Hence, for every \olup{structure}
$\Struct M$ with $\Sat{M}{\Gamma}$, $\Sat{M}{A}$, so $\Gamma \Entails
A$.
\end{proof}

\begin{prob}
\begin{enumerate}
\item Show that $\Gamma \Entails \bot$ iff $\Gamma$ is unsatisfiable.
\item Show that $\Gamma \cup \{A\} \Entails \bot$ iff $\Gamma \Entails \lnot A$.
\item Suppose $c$ does not occur in $A$ or $\Gamma$.  Show that
  $\Gamma \Entails \lforall[x][A]$ iff $\Gamma \Entails
  \Subst{A}{c}{x}$.
\end{enumerate}
\end{prob}

\begin{prop}
If $\Gamma \subseteq \Gamma'$ and $\Gamma \Entails A$, then $\Gamma'
\Entails A$.
\end{prop}

\begin{proof}
Suppose that $\Gamma \subseteq \Gamma'$ and $\Gamma \Entails A$. Let
$\Struct M$ be a structure such that $\Sat{M}{\Gamma'}$; then $\Sat{M}{\Gamma}$,
and since $\Gamma \Entails A$, we get that $\Sat{M}{A}$. Hence,
whenever $\Sat{M}{\Gamma'}$, $\Sat{M}{A}$, so $\Gamma' \Entails A$.
\end{proof}


\begin{thm}[Semantic Deduction Theorem]
\ollabel{thm:sem-deduction}
$\Gamma \cup \{A\} \Entails B$ iff $\Gamma \Entails A \lif B$.
\end{thm}

\begin{proof}
For the forward direction, let $\Gamma \cup \{ A \} \Entails B$ and
let $\Struct M$ be \olup{structure} so that $\Sat{M}{\Gamma}$. If
$\Sat{M}{A}$, then $\Sat{M}{\Gamma \cup \{ A \} }$, so since $\Gamma
\cup \{ A \}$ entails $B$, we get $\Sat{M}{B}$. Therefore,
$\Sat{M}{A \lif B}$, so $\Gamma \Entails A \lif B$.

For the reverse direction, let $\Gamma \Entails A \lif B$ and
$\Struct M$ be \olup{structure} so that $\Sat{M}{\Gamma \cup \{ A
  \}}$. Then $\Sat{M}{\Gamma}$, so $\Sat{M}{A \lif B}$, and since
$\Sat{M}{A}$, $\Sat{M}{B}$. Hence, whenever $\Sat{M}{\Gamma \cup \{
  A \} }$, $\Sat{M}{B}$, so $\Gamma \cup \{ A \} \Entails B$.
\end{proof}

\begin{prop}\ollabel{prop:quant-terms}
  Let $\Struct{M}$ be \olup{structure}, and $A(x)$ \olup{formula} with
  one free variable~$x$, and $t$~a closed term. Then:
  \begin{tagenumerate}{prvEx,prvAll}
    \tagitem{prvEx}{$A(t) \Entails \lexists[x][A(x)]$}{}
    \tagitem{prvAll}{$\lforall[x][A(x)] \Entails A(t)$}{}
  \end{tagenumerate}
\end{prop}

\begin{proof}
  \begin{tagenumerate}{prvEx,prvAll}
    \tagitem{prvEx}{%
      \iftag{probEx}{Exercise.}{Suppose $\Sat{M}{A(t)}$. Let $s$ be a
        variable assignment with $s(x) = \Value{t}{M}$. Then
        $\Sat{M}{A(t)}[s]$ since $A(t)$ is \olup{sentence}. By
        \olref[ext]{prop:ext-formulas}, $\Sat{M}{A(x)}[s]$. By
        \olref[ass]{prop:sat-quant},
        $\Sat{M}{\lexists[x][A(x)]}$.}}{}
    \tagitem{prvAll}{%
      \iftag{probAll}{Exercise.}{Suppose
        $\Sat{M}{\lforall[x][A(x)]}$. Let $s$ be a variable
        assignment with $s(x) = \Value{t}{M}$. By
        \olref[ass]{prop:sat-quant}, $\Sat{M}{A(x)}[s]$.  By
        \olref[ext]{prop:ext-formulas}, $\Sat{M}{A(t)}[s]$. By
        \olref[ass]{prop:sentence-sat-true}, $\Sat{M}{A(t)}$ since
        $A(t)$ is \olup{sentence}.}}{}
  \end{tagenumerate}
\end{proof}

\begin{probtag}{probEx,probAll}
  Complete the proof of \olref[fol][syn][sem]{prop:quant-terms}.
\end{probtag} 


% ==== OLP file: content/first-order-logic/syntax-and-semantics/semantics.tex ====
% Part: first-order-logic
% Chapter: semantics



\olchapter{fol}{sem}{Semantics of First-Order Logic}

\olimport{intro-semantics}

\olimport{structures}

\olimport{covered-structures}

\olimport{satisfaction}

\olimport{assignments}

\olimport{extensionality}

\olimport{semantic-notions}

\OLEndChapterHook


% ==== OLP file: content/first-order-logic/syntax-and-semantics/structures.tex ====
% Part: first-order-logic
% Chapter: syntax-and-semantics
% Section: Structures



\olfileid{fol}{syn}{str}

\olsection{\printtoken{P}{structure} for First-order Languages}

\begin{explain}
First-order languages are, by themselves, \emph{uninterpreted:} the
\olup{constant}s, \olup{function}s, and \olup{predicate}s have no specific
meaning attached to them.  Meanings are given by specifying
\article{structure} \emph{\olup{structure}}. It specifies the
\emph{domain}, i.e., the objects which the \olup{constant}s pick out, the
\olup{function}s operate on, and the quantifiers range over. In addition,
it specifies which \olup{constant}s pick out which objects, how
\olup{function} maps objects to objects, and which objects the
\olup{predicate}s apply to.  !!^{structure}s are the basis for
\emph{semantic} notions in logic, e.g., the notion of consequence,
validity, satisfiability. They are variously called ``structures,''
``interpretations,'' or ``models'' in the literature.
\end{explain}

\begin{defn}[!!^{structure}s]
\Article{structure} \emph{\olup{structure}}~$\Struct M$, for a language
$\Lang{L}$ of first-order logic consists of the following elements:
\begin{enumerate}
\item \emph{Domain:} a non-empty set, $\Domain M$
\item \emph{Interpretation of \olup{constant}s:} for each \olup{constant}~$c$ of
  $\Lang{L}$, \olup{element} $\Assign{c}{M} \in \Domain M$
\item \emph{Interpretation of \olup{predicate}s:} for each $n$-place
  \olup{predicate}~$R$ of $\Lang{L}$ (other than $\eq$), an $n$-place
  relation $\Assign{R}{M} \subseteq \Domain{M}^n$
\item \emph{Interpretation of \olup{function}s:} for each $n$-place
  \olup{function}~$f$ of $\Lang{L}$, an $n$-place function $\Assign{f}{M}
  \colon \Domain{M}^n \to \Domain{M}$
\end{enumerate}
\end{defn}

\begin{ex}
!!^a{structure}~$\Struct M$ for the language of arithmetic consists of a
  set, an element of $\Domain M$, $\Assign{\Obj 0}{M}$, as
  interpretation of the \olup{constant}~$\Obj 0$, a one-place function
  $\Assign{\Obj \prime}{M} \colon \Domain{M} \to \Domain M$, two
  two-place functions $\Assign{\Obj +}{M}$ and $\Assign{\Obj
    \times}{M}$, both $\Domain M^2 \to \Domain M$, and a two-place
  relation $\Assign{\Obj <}{M} \subseteq \Domain{M}^2$.

An obvious example of such a structure is the following:
\begin{enumerate}
\item $\Domain N = \Nat$
\item $\Assign{\Obj 0}{N} = 0$
\item $\Assign{\Obj \prime}{N}(n) = n + 1$ for all $n \in \Nat$
\item $\Assign{\Obj +}{N}(n, m) = n + m$ for all $n, m \in \Nat$
\item $\Assign{\Obj \times}{N}(n, m) = n\cdot m$ for all $n, m \in \Nat$
\item $\Assign{\Obj <}{N} = \Setabs{\tuple{n, m}}{n \in \Nat, m \in
  \Nat, n < m}$
\end{enumerate}
The structure~$\Struct N$ for $\Lang L_A$ so defined is called the
\emph{standard model of arithmetic}, because it interprets the
non-logical constants of~$\Lang L_A$ exactly how you would expect.

However, there are many other possible \olup{structure}s for~$\Lang
L_A$. For instance, we might take as the domain the set~$\Int$ of
integers instead of~$\Nat$, and define the interpretations of $\Obj
0$, $\Obj \prime$, $\Obj +$, $\Obj \times$, $\Obj <$ accordingly.  But
we can also define structures for~$\Lang L_A$ which have nothing even
remotely to do with numbers.
\end{ex}

\begin{ex}
A structure~$\Struct M$ for the language~$\Lang L_Z$ of set theory requires
just a set and a single-two place relation. So technically, e.g., the
set of people plus the relation ``$x$ is older than $y$'' could be
used as \olup{structure} for $\Lang L_Z$, as well as $\Nat$ together
with $n \ge m$ for $n, m \in \Nat$.

A particularly interesting \olup{structure} for $\Lang L_Z$ in which the
\olup{element}s of the domain are actually sets, and the interpretation
of $\Obj \in$ actually is the relation ``$x$ is \olup{element} of~$y$''
is the \olup{structure}~$\Struct{{HF}}$ of \emph{hereditarily finite sets}:
\begin{enumerate}
\item $\Domain{{{HF}}} = \emptyset \cup \Pow{\emptyset} \cup
  \Pow{\Pow{\emptyset}} \cup \Pow{\Pow{\Pow{\emptyset}}} \cup \dots$;
\item $\Assign{\Obj \in}{{{HF}}} = \Setabs{\tuple{x, y}}{x, y \in
  \Domain{{{HF}}}, x \in y}$.
\end{enumerate}
\end{ex}

\begin{digress}
The stipulations we make as to what counts as \olup{structure} impact
our logic. For example, the choice to prevent empty domains ensures,
given the usual account of satisfaction (or truth) for quantified
sentences, that $\lexists[x][(A(x) \lor \lnot A(x))]$ is
valid---that is, a logical truth. And the stipulation that all
\olup{constant}s must refer to an object in the domain ensures that the
existential generalization is a sound pattern of inference: $A(a)$,
therefore $\lexists[x][A(x)]$. If we allowed names to refer outside
the domain, or to not refer, then we would be on our way to a
\emph{free logic}, in which existential generalization requires an
additional premise: $A(a)$ and $\lexists[x][\eq[x][a]]$, therefore
$\lexists[x][A(x)]$.
\end{digress}


% ==== OLP file: content/first-order-logic/syntax-and-semantics/subformulas.tex ====
% Part: first-order-logic
% Chapter: syntax-and-semantics
% Section: subformulas



\olfileid{fol}{syn}{sbf}

\olsection{\printtoken{P}{subformula}}

\begin{explain}
It is often useful to talk about the \olup{formula}s that ``make up'' a
given \olup{formula}.  We call these its \emph{\olup{subformula}s}.  Any
\olup{formula} counts as \olup{subformula} of itself; a subformula of $A$
other than $A$ itself is a \emph{proper \olup{subformula}}.
\end{explain}

\begin{defn}[Immediate !!^{subformula}]
If $A$ is \olup{formula}, the \emph{immediate \olup{subformula}s}
of $A$ are defined inductively as follows:
\begin{enumerate}
\item Atomic \olup{formula}s have no immediate \olup{subformula}s.

\tagitem{prvNot}{\indcase{A}{\lnot B}{The only immediate
    \olup{subformula} of $\indfrm$ is~$B$.}}{}

\item \indcase{A}{(B \ast C)}{The immediate \olup{subformula}s of
  $\indfrm$ are $B$ and $C$ ($\ast$ is any one of the two-place
  connectives).}

\tagitem{prvAll}{\indcase{A}{\lforall[x][B]}{The only immediate
    \olup{subformula} of $\indfrm$ is~$B$.}}{}

\tagitem{prvEx}{\indcase{A}{\lexists[x][B]}{The only immediate
    \olup{subformula} of $\indfrm$ is~$B$.}}{}
\end{enumerate}
\end{defn}

\begin{defn}[Proper !!^{subformula}]
If $A$ is \olup{formula}, the \emph{proper \olup{subformula}s}
of $A$ are defined recursively as follows:
\begin{enumerate}
\item Atomic \olup{formula}s have no proper \olup{subformula}s.

\tagitem{prvNot}{\indcase{A}{\lnot B}{The proper \olup{subformula}s of
    $\indfrm$ are~$B$ together with all proper \olup{subformula}s
    of~$B$.}}{}

\item \indcase{A}{(B \ast C)}{The proper \olup{subformula}s of
  $\indfrm$ are $B$, $C$, together with all proper \olup{subformula}s
  of $B$ and those of~$C$.}

\tagitem{prvAll}{\indcase{A}{\lforall[x][B]}{The proper
    \olup{subformula}s of $\indfrm$ are~$B$ together with all proper
    \olup{subformula}s of~$B$.}}{}

\tagitem{prvEx}{\indcase{A}{\lexists[x][B]}{The proper
    \olup{subformula}s of $\indfrm$ are~$B$ together with all proper
    \olup{subformula}s of~$B$.}}{}
\end{enumerate}
\end{defn}

\begin{defn}[!!^{subformula}]
The \olup{subformula}s of $A$ are $A$ itself together with all its
proper \olup{subformula}s.
\end{defn}

\begin{explain}
Note the subtle difference in how we have defined immediate
\olup{subformula}s and proper \olup{subformula}s.  In the first case, we
have directly defined the immediate \olup{subformula}s of a formula~$A$
for each possible form of~$A$.  It is an explicit definition by
cases, and the cases mirror the inductive definition of the set of
\olup{formula}s.  In the second case, we have also mirrored the way the
set of all \olup{formula}s is defined, but in each case we have also
included the proper \olup{subformula}s of the smaller \olup{formula}s $B$,
$C$ in addition to these \olup{formula}s themselves.  This makes the
definition \emph{recursive}.  In general, a definition of a function
on an inductively defined set (in our case, \olup{formula}s) is recursive
if the cases in the definition of the function make use of
the function itself. To be well defined, we must make sure, however,
that we only ever use the values of the function for arguments that
come ``before'' the one we are defining---in our case, when defining
``proper \olup{subformula}'' for $(B \ast C)$ we only use the proper
\olup{subformula}s of the ``earlier'' \olup{formula}s $B$ and $C$.
\end{explain}

\begin{prop}
\ollabel{prop:subfrm-trans}
Suppose $B$ is a subformula of $A$ and $C$ is a subformula of $B$.
Then $C$ is a subformula of $A$. In other words, the subformula
relation is transitive.
\end{prop}

\begin{prob}
Prove \olref[fol][syn][sbf]{prop:subfrm-trans}.
\end{prob}

\begin{prop}
\ollabel{prop:count-subfrms}
Suppose $A$ is a formula with $n$ connectives and quantifiers.
Then $A$ has at most $2n+1$ subformulas.
\end{prop}

\begin{prob}
Prove \olref[fol][syn][sbf]{prop:count-subfrms}.
\end{prob}


% ==== OLP file: content/first-order-logic/syntax-and-semantics/substitution.tex ====
% Part: first-order-logic
% Chapter: syntax-and-semantics
% Section: substitution



\olfileid{fol}{syn}{sub}

\olsection{Substitution}

\begin{defn}[Substitution in a term]
We define $\Subst{s}{t}{x}$, the result of \emph{substituting} $t$
for every occurrence of~$x$ in $s$, recursively:
\begin{enumerate}
\item \indcase{s}{c}{$\Subst{\indfrm}{t}{x}$ is just $s$.}

\item \indcase{s}{y}{$\Subst{\indfrm}{t}{x}$ is also just~$s$,
  provided $y$ is a variable and $y \not\ident x$.}

\item \indcase{s}{x}{$\Subst{\indfrm}{t}{x}$ is~$t$.}

\item\indcase{s}{\Atom{f}{t_1, \dots, t_n}}{$\Subst{\indfrmp}{t}{x}$ is
  $\Atom{f}{\Subst{t_1}{t}{x}, \dots, \Subst{t_n}{t}{x}}$.}
\end{enumerate}
\end{defn}

\begin{defn}
A term~$t$ is \emph{\olup{free for}} $x$ in $A$ if none of the free
occurrences of~$x$ in $A$ occur in the scope of a quantifier that
binds a variable in~$t$.
\end{defn}

\begin{ex} ~
\begin{enumerate}
\item $\Obj v_8$ is free for $\Obj v_1$ in $\lexists[\Obj
  v_3]\Atom{\Obj A^2_4}{\Obj v_3,\Obj v_1}$

\item $\Obj f^2_1(\Obj v_1, \Obj v_2)$ is \emph{not} free for $\Obj
  v_0$ in $\lforall[\Obj v_2]\Atom{\Obj A^2_4}{\Obj v_0,\Obj v_2}$
\end{enumerate}
\end{ex}

\begin{defn}[Substitution in \olup{formula}]
If $A$ is \olup{formula}, $x$~is \olup{variable}, and $t$~is a term
\olup{free for}~$x$ in~$A$, then $\Subst{A}{t}{x}$ is the result of
substituting $t$ for all free occurrences of~$x$ in~$A$.
\begin{enumerate}
\tagitem{prvFalse}{\indcase{A}{\lfalse}{$\Subst{\indfrm}{t}{x}$ is
    $\lfalse$.}}{}

\tagitem{prvTrue}{\indcase{A}{\ltrue}{$\Subst{\indfrm}{t}{x}$ is
    $\ltrue$.}}{}

\item \indcase{A}{\Atom{P}{t_1,\dots,
    t_n}}{$\Subst{\indfrm}{t}{x}$ is $\Atom{P}{\Subst{t_1}{t}{x},
    \dots, \Subst{t_n}{t}{x}}$.}

\item \indcase{A}{\eq[t_1][t_2]}{$\Subst{\indfrmp}{t}{x}$ is
  $\Subst{t_1}{t}{x} = \Subst{t_2}{t}{x}$.}

\tagitem{prvNot}{\indcase{A}{\lnot B}{$\Subst{\indfrmp}{t}{x}$ is
    $\lnot \Subst{B}{t}{x}$.}}{}

\tagitem{prvAnd}{\indcase{A}{(B \land
    C)}{$\Subst{\indfrmp}{t}{x}$ is $(\Subst{B}{t}{x} \land
    \Subst{C}{t}{x})$.}}{}

\tagitem{prvOr}{\indcase{A}{(B \lor
    C)}{$\Subst{\indfrmp}{t}{x}$ is $(\Subst{B}{t}{x} \lor
    \Subst{C}{t}{x})$.}}{}

\tagitem{prvIf}{\indcase{A}{(B \lif
    C)}{$\Subst{\indfrmp}{t}{x}$ is $(\Subst{B}{t}{x} \lif
    \Subst{C}{t}{x})$.}}{}

\tagitem{prvIff}{\indcase{A}{(B \liff
    C)}{$\Subst{\indfrmp}{t}{x}$ is $(\Subst{B}{t}{x} \liff
    \Subst{C}{t}{x})$.}}{}

\tagitem{prvAll}{
\indcase{A}{\lforall[y][B]}{$\Subst{\indfrmp}{t}{x}$
    is $\lforall[y][\Subst{B}{t}{x}]$, provided $y$ is a variable
    other than $x$; otherwise $\Subst{\indfrmp}{t}{x}$
    is just $\indfrm$.}}{}

\tagitem{prvEx}{
\indcase{A}{\lexists[y][B]}{$\Subst{\indfrmp}{t}{x}$
    is $\lexists[y][\Subst{B}{t}{x}]$, provided $y$ is a variable
    other than $x$; otherwise $\Subst{\indfrmp}{t}{x}$
    is just $\indfrm$.}}{}
\end{enumerate}
\end{defn}

\begin{explain}
Note that substitution may be vacuous: If $x$ does not occur in $A$
at all, then $\Subst{A}{t}{x}$ is just~$A$.

The restriction that $t$ must be \olup{free for}~$x$ in~$A$ is necessary to
exclude cases like the following. If $A \ident \lexists[y][x < y]$
and $t \ident y$, then $\Subst{A}{t}{x}$ would be $\lexists[y][y <
  y]$. In this case the free variable $y$ is ``captured'' by the
quantifier $\lexists[y]$ upon substitution, and that is undesirable.
For instance, we would like it to be the case that whenever
$\lforall[x][B]$ holds, so does $\Subst{B}{t}{x}$. But consider
$\lforall[x][\lexists[y][x < y]]$ (here $B$ is $\lexists[y][x <
  y]$). It is a sentence that is true about, e.g., the natural numbers:
for every number~$x$ there is a number~$y$ greater than it. If we
allowed $y$ as a possible substitution for~$x$, we would end up with
$\Subst{B}{y}{x} \ident \lexists[y][y < y]$, which is false. We
prevent this by requiring that none of the free variables in~$t$ would
end up being bound by a quantifier in~$A$.
\end{explain}

We often use the following convention to avoid cumbersome notation: If
$A$ is \olup{formula} which may contain the \olup{variable}~$x$ free, we
also write~$A(x)$ to indicate this. When it is clear which $A$
and~$x$ we have in mind, and $t$ is a term (assumed to be free for $x$
in $A(x)$), then we write $A(t)$ as short for $\Subst{A}{t}{x}$. So
for instance, we might say, ``we call $A(t)$ an instance
of~$\lforall[x][A(x)]$.'' By this we mean that if $A$~is any
\olup{formula}, $x$~\olup{variable}, and $t$~a term that's free for~$x$
in~$A$, then $\Subst{A}{t}{x}$ is an instance of~$\lforall[x][A]$.


% ==== OLP file: content/first-order-logic/syntax-and-semantics/syntax-and-semantics.tex ====
% Part: first-order-logic
% Chapter: syntax-and-semantics



\olchapter{fol}{syn}{Syntax and Semantics}

\olimport{introduction}

\olimport{first-order-languages}

\olimport{terms-formulas}

\olimport{unique-readability}

\olimport{main-operator}

\olimport{subformulas}

\olimport{formation-sequences}

\olimport{free-vars-sentences}

\olimport{substitution}

\olimport{structures}

\olimport{covered-structures}

\olimport{satisfaction}

\olimport{assignments}

\olimport{extensionality}

\olimport{semantic-notions}

\OLEndChapterHook


% ==== OLP file: content/first-order-logic/syntax-and-semantics/syntax.tex ====
% Part: first-order-logic
% Chapter: syntax



\olchapter{fol}{syn}{Syntax of First-Order Logic}

\olimport{intro-syntax}

\olimport{first-order-languages}

\olimport{terms-formulas}

\olimport{unique-readability}

\olimport{main-operator}

\olimport{subformulas}

\olimport{formation-sequences}

\olimport{free-vars-sentences}

\olimport{substitution}

\OLEndChapterHook


% ==== OLP file: content/first-order-logic/syntax-and-semantics/terms-formulas.tex ====
% Part: first-order-logic
% Chapter: syntax-and-semantics
% Section: terms-formulas



\olfileid{fol}{syn}{frm}

\olsection{Terms and \printtoken{P}{formula}}

Once a first-order language~$\Lang L$ is given, we can define
expressions built up from the basic vocabulary of~$\Lang L$.  These
include in particular \emph{terms} and \emph{\olup{formula}s}.

\begin{defn}[Terms]
\ollabel{defn:terms}
The set of \emph{terms}~$\Trm[L]$ of~$\Lang L$ is
defined inductively by:
\begin{enumerate}
\item Every \olup{variable} is a term.
\item Every \olup{constant} of~$\Lang L$ is a term.
\item If $f$ is an $n$-place \olup{function} and $t_1$, \dots, $t_n$
  are terms, then $\Atom{f}{t_1, \ldots, t_n}$ is a term.
\tagitem{limitClause}{
  Nothing else is a term.}{}
\end{enumerate}
A term containing no \olup{variable}s is a \emph{closed term}.
\end{defn}

\begin{explain}
The \olup{constant}s appear in our specification of the language and the
terms as a separate category of symbols, but they could instead have
been included as zero-place \olup{function}s.  We could then do without
the second clause in the definition of terms. We just have to
understand $\Atom{f}{t_1, \ldots, t_n}$ as just $f$ by itself if $n =
0$.
\end{explain}

\begin{defn}[Formulas]
\ollabel{defn:formulas}
The set of \emph{\olup{formula}s}~$\Frm[L]$ of the language~$\Lang L$
is defined inductively as follows:
\begin{enumerate}
\tagitem{prvFalse}{$\lfalse$ is an atomic \olup{formula}.}{}

\tagitem{prvTrue}{$\ltrue$ is an atomic \olup{formula}.}{}

\item If $R$ is an $n$-place \olup{predicate} of~$\Lang L$ and $t_1$, \dots,
  $t_n$ are terms of~$\Lang L$, then $\Atom{R}{t_1,\ldots, t_n}$ is an
  atomic \olup{formula}.

\item If $t_1$ and $t_2$ are terms of~$\Lang L$, then $\Atom{\eq}{t_1, t_2}$
  is an atomic \olup{formula}.

\tagitem{prvNot}{If $A$ is \olup{formula}, then $\lnot A$ is
  \olup{formula}.}{}

\tagitem{prvAnd}{If $A$ and $B$ are \olup{formula}s, then $(A \land
  B)$ is \olup{formula}.}{}

\tagitem{prvOr}{If $A$ and $B$ are \olup{formula}s, then $(A \lor B)$
  is \olup{formula}.}{}

\tagitem{prvIf}{If $A$ and $B$ are \olup{formula}s, then $(A \lif B)$
  is \olup{formula}.}{}

\tagitem{prvIff}{If $A$ and $B$ are \olup{formula}s, then $(A \liff B)$
  is \olup{formula}.}{}

\tagitem{prvAll}{If $A$ is \olup{formula} and $x$ is \olup{variable},
  then $\lforall[x][A]$ is \olup{formula}.}{}

\tagitem{prvEx}{If $A$ is \olup{formula} and $x$ is \olup{variable},
  then $\lexists[x][A]$ is \olup{formula}.}{}

\tagitem{limitClause}{Nothing else is \olup{formula}.}{}
\end{enumerate}
\end{defn}

\begin{explain}
The definitions of the set of terms and that of \olup{formula}s are
\emph{inductive definitions}.  Essentially, we construct the set of
\olup{formula}s in infinitely many stages.  In the initial stage, we
pronounce all atomic formulas to be formulas; this corresponds to the
first few cases of the definition, i.e., the cases for
\iftag{prvTrue}{$\ltrue$, }{}%
\iftag{prvFalse}{$\lfalse$, }{}%
$\Atom{R}{t_1,\dots,t_n}$ and $\Atom{\eq}{t_1,t_2}$. ``Atomic
\olup{formula}'' thus means any \olup{formula} of this form.

The other cases of the definition give rules for constructing new
\olup{formula}s out of \olup{formula}s already constructed.  At the second
stage, we can use them to construct \olup{formula}s out of atomic
\olup{formula}s.  At the third stage, we construct new formulas from the
atomic formulas and those obtained in the second stage, and so on.  A
\olup{formula} is anything that is eventually constructed at such a
stage, and nothing else.
\end{explain}

By convention, we write $\eq$ between its arguments and leave out the
parentheses: $\eq[t_1][t_2]$ is an abbreviation for
$\Atom{\eq}{t_1,t_2}$.  Moreover, $\lnot \Atom{\eq}{t_1,t_2}$ is
abbreviated as $\eq/[t_1][t_2]$. When writing a formula $(B \ast C)$
constructed from $B$, $C$ using a two-place connective~$\ast$, we
will often leave out the outermost pair of parentheses and write
simply~$B \ast C$.

\begin{intro}
Some logic texts require that the \olup{variable}~$x$ must occur in~$A$
in order for
\iftag{prvEx}{$\lexists[x][A]$ }{}%
\iftag{notprvEx,notprvAll}{}{and }%
\iftag{prvAll}{$\lforall[x][A]$ }{}%
to count as
\iftag{notprvEx,notprvAll}{\olup{formula}}{\olup{formula}s}.
Nothing bad happens if you don't require this, and it makes things
easier.
\end{intro}

\begin{tagblock}{defNot,defOr,defAnd,defIf,defIff,defTrue,defFalse,defEx,defAll}
\begin{defn}
Formulas constructed using the defined operators are to be understood
as follows:

\begin{tagenumerate}{defTrue,defFalse,defNot,defOr,defAnd,defIf,defIff,defEx,defAll}

\tagitem{defTrue}{$\ltrue$ abbreviates
  \iftag{prvFalse}{$\lnot\lfalse$}{$(A \lor \lnot A)$ for some
    fixed atomic \olup{formula}~$A$}.}{}

\tagitem{defFalse}{$\lfalse$ abbreviates
  \iftag{prvTrue}{$\lnot\ltrue$}{$(A \land \lnot A)$ for some
    fixed atomic \olup{formula}~$A$}.}{}

\tagitem{defNot}{$\lnot A$ abbreviates $A \lif \lfalse$.}{}

\tagitem{defOr}{$A \lor B$ abbreviates
  \iftag{prvAnd}{$\lnot(\lnot A \land \lnot B)$}{$\lnot A \lif
    B$}.}{}

\tagitem{defAnd}{$A \land B$ abbreviates
  \iftag{prvOr}{$\lnot(\lnot A \lor \lnot B)$}{$\lnot (A \lif
    \lnot B)$}.}{}

\tagitem{defIf}{$A \lif B$ abbreviates
  \iftag{prvOr}{$\lnot A \lor B)$}{$\lnot (A \land \lnot B)$}.}{}

\tagitem{defIff}{$A \liff B$ abbreviates $(A \lif B) \land (B
  \lif A)$.}{}

\tagitem{defAll}{$\lforall[x][A]$ abbreviates $\lnot\lexists[x][\lnot A]$.}{}

\tagitem{defEx}{$\lexists[x][A]$ abbreviates $\lnot\lforall[x][\lnot A]$.}{}
\end{tagenumerate}
\end{defn}
\end{tagblock}

If we work in a language for a specific application, we will often
write two-place \olup{predicate}s and \olup{function}s between the
respective terms, e.g., $t_1 < t_2$ and $(t_1 + t_2)$ in the language
of arithmetic and $t_1 \in t_2$ in the language of set theory.  The
successor function in the language of arithmetic is even written
conventionally \emph{after} its argument:~$t'$.  Officially, however,
these are just conventional abbreviations for $\Atom{\Obj A^2_0}{t_1, t_2}$,
$\Obj f^2_0(t_1, t_2)$, $\Atom{\Obj A^2_0}{t_1, t_2}$ and $f^1_0(t)$,
respectively.

\begin{defn}[Syntactic identity]
The symbol $\ident$ expresses syntactic identity between strings of
symbols, i.e., $A \ident B$ iff $A$ and $B$ are strings of symbols
of the same length and which contain the same symbol in each place.
\end{defn}

The $\ident$ symbol may be flanked by strings obtained by
concatenation, e.g., $A \ident (B \lor C)$ means: the string of
symbols~$A$ is the same string as the one obtained by concatenating
an opening parenthesis, the string $B$, the $\lor$ symbol, the
string~$C$, and a closing parenthesis, in this order. If this is the
case, then we know that the first symbol of~$A$ is an opening
parenthesis, $A$ contains $B$ as a substring (starting at the second
symbol), that substring is followed by~$\lor$, etc.

As terms and \olup{formula}s are built up from basic elements via inductive
definitions, we can use the following induction principles to prove
things about them.

\begin{lem}[\emph{Principle of induction on terms}]
    \ollabel{lem:trmind}
    Let $\Lang L$ be a first-order language.
    If some property~$P$ is such that
    %
    \begin{enumerate}
        \item it holds for every \olup{variable}~$v$,
        %
        \item it holds for every \olup{constant}~$a$ of~$\Lang L$, and
        %
        \item it holds for $f(t_1,\dotsc,t_n)$ whenever it holds for
        $t_1$,~\dots, $t_n$ and $f$~is an $n$-place
            \olup{function} of~$\Lang L$
    \end{enumerate}
    (assuming $t_1$,~\dots, $t_n$ are terms of~$\Lang{L}$),
    then $P$ holds for every term in~$\Trm[L]$.
\end{lem}

\begin{prob}
    Prove \olref[fol][syn][frm]{lem:trmind}.
\end{prob}

\begin{lem}[\emph{Principle of induction on \olup{formula}s}]
    \ollabel{thm:frmind}
    Let $\Lang L$ be a first-order language.
    If some property~$P$ holds for all the atomic \olup{formula}s
    and is such that
    %
    \begin{enumerate}
        \tagitem{prvNot}{it holds for $\lnot A$ whenever it
            holds for~$A$;}{}
        \tagitem{prvAnd}{it holds for $(A \land B)$
            whenever it holds for $A$ and~$B$;}{}
        \tagitem{prvOr}{it holds for $(A \lor B)$
            whenever it holds for $A$ and~$B$;}{}
        \tagitem{prvIf}{it holds for $(A \lif B)$
            whenever it holds for $A$ and~$B$;}{}
        \tagitem{prvIff}{it holds for $(A \liff B)$
            whenever it holds for $A$ and~$B$;}{}
        \tagitem{prvEx}{it holds for $\lexists[x][A]$
            whenever it holds for~$A$;}{}
        \tagitem{prvAll}{it holds for $\lforall[x][A]$
            whenever it holds for~$A$;}{}
  \end{enumerate}
  (assuming $A$ and $B$ are \olup{formula}s of~$\Lang{L}$),
  then $P$ holds for all formulas in~$\Frm[L]$.
\end{lem}


% ==== OLP file: content/first-order-logic/syntax-and-semantics/unique-readability.tex ====
% Part: first-order-logic
% Chapter: syntax-and-semantics
% Section: unique-readability



\olfileid{fol}{syn}{unq}

\olsection{Unique Readability}

\begin{explain}
The way we defined \olup{formula}s guarantees that every \olup{formula} has
a \emph{unique reading}, i.e., there is essentially only one way of
constructing it according to our formation rules for \olup{formula}s and
only one way of ``interpreting'' it.  If this were not so, we would
have ambiguous \olup{formula}s, i.e., \olup{formula}s that have more than
one reading or intepretation---and that is clearly something we want
to avoid.  But more importantly, without this property, most of the
definitions and proofs we are going to give will not go through.

Perhaps the best way to make this clear is to see what would happen if
we had given bad rules for forming \olup{formula}s that would not
guarantee unique readability.  For instance, we could have forgotten
the parentheses in the formation rules for connectives, e.g., we might
have allowed this:
\begin{quote}
If $A$ and $B$ are \olup{formula}s, then so is $A \lif B$.
\end{quote}
Starting from an atomic formula $D$, this would allow us to form $D
\lif D$. From this, together with $D$, we would get $D \lif D
\lif D$.  But there are two ways to do this:
\begin{enumerate}
\item We take $D$ to be $A$ and $D \lif D$ to be $B$.
\item We take $A$ to be $D \lif D$ and $B$ is~$D$.  
\end{enumerate}
Correspondingly, there are two ways to
``read'' the \olup{formula}~$D \lif D \lif D$. It is of the form $B
\lif C$ where $B$ is $D$ and $C$ is $D \lif D$, but \emph{it is
  also} of the form $B \lif C$ with $B$ being $D \lif D$ and $C$
being~$D$.

If this happens, our definitions will not always work. For instance,
when we define the \olup{main operator} of a formula, we say: in a
formula of the form $B \lif C$, the \olup{main operator} is the
indicated occurrence of~$\lif$. But if we can match the formula $D
\lif D \lif D$ with $B \lif C$ in the two different ways mentioned
above, then in one case we get the first occurrence of $\lif$ as the
\olup{main operator}, and in the second case the second occurrence.  But
we intend the \olup{main operator} to be a \emph{function} of the
\olup{formula}, i.e., every \olup{formula} must have exactly one \olup{main
  operator} occurrence.
\end{explain}

\begin{lem}
The number of left and right parentheses in \olup{formula}~$A$ are
equal.
\end{lem}

\begin{proof}
We prove this by induction on the way $A$ is constructed. This
requires two things: (a) We have to prove first that all atomic
formulas have the property in question (the induction basis). (b) Then
we have to prove that when we construct new formulas out of given
formulas, the new formulas have the property provided the old ones do.

Let $l(A)$ be the number of left parentheses, and $r(A)$ the number
of right parentheses in~$A$, and $l(t)$ and $r(t)$ similarly the
number of left and right parentheses in a term~$t$.

\begin{prob}
    Prove that for any term~$t$, $l(t) = r(t)$.
\end{prob}

\begin{enumerate}
\tagitem{prvFalse}{\indcase{A}{\lfalse}{$\indfrm$ has $0$ left and $0$
    right parentheses.}}{}

\tagitem{prvTrue}{\indcase{A}{\ltrue}{$\indfrm$ has $0$ left and $0$
    right parentheses.}}{}

\item \indcase{A}{\Atom{R}{t_1,\dots,t_n}}{$l(\indfrm) = 1 + l(t_1) +
  \dots + l(t_n) = 1 + r(t_1) + \dots + r(t_n) = r(\indfrm)$. Here we
  make use of the fact, left as an exercise, that $l(t) = r(t)$ for
  any term~$t$.}

\item \indcase{A}{\eq[t_1][t_2]}{$l(\indfrm) = l(t_1) + l(t_2) =
  r(t_1) + r(t_2) = r(\indfrm)$.}

\tagitem{prvNot}{\indcase{A}{\lnot B}{By induction hypothesis,
    $l(B) = r(B)$. Thus $l(\indfrm) = l(B) = r(B) =
    r(\indfrm)$.}}{}

\item \indcase{A}{(B \ast C)}{By induction hypothesis, $l(B) =
  r(B)$ and $l(C) = r(C)$.  Thus $l(\indfrm) = 1 + l(B) + l(C) =
  1 + r(B) + r(C) = r(\indfrm)$.}

\tagitem{prvAll}{\indcase{A}{\lforall[x][B]}{By induction
    hypothesis, $l(B) = r(B)$. Thus, $l(\indfrm) = l(B) = r(B) =
    r(\indfrm)$.}}{}

% Print case for \lexists if prvEx and not prvAll
\tagitem{prvAll,notprvEx}{}{\indcase{A}{\lexists[x][B]}{By induction
    hypothesis, $l(B) = r(B)$. Thus, $l(\indfrm) = l(B) = r(B) =
    r(\indfrm)$.}}

% Just say similarly if prvEx and prvAll
\tagitem{notprvAll,notprvEx}{}{\indcase{A}{\lexists[x][B]}{Similarly.}}
\end{enumerate}
\end{proof}

\begin{defn}[Proper prefix]
A string of symbols $B$ is a \emph{proper prefix} of a string of symbols~$A$ if
concatenating $B$ and a non-empty string of symbols yields~$A$.
\end{defn}

\begin{lem}\ollabel{lem:no-prefix}
If $A$ is \olup{formula}, and $B$ is a proper prefix of $A$, then
$B$ is not \olup{formula}.
\end{lem}

\begin{proof}
Exercise.
\end{proof}

\begin{prob}
Prove \olref[fol][syn][unq]{lem:no-prefix}.
\end{prob}

\begin{prop}
\ollabel{prop:unique-atomic}
If $A$ is an atomic \olup{formula}, then it satisfies one, and only one
of the following conditions.
\begin{enumerate}
\tagitem{prvFalse}{$A \ident \lfalse$.}{}
\tagitem{prvTrue}{$A \ident \ltrue$.}{}
\item $A \ident \Atom{R}{t_1,\dots,t_n}$ where $R$ is an $n$-place
  \olup{predicate}, $t_1$, \dots, $t_n$ are terms, and each of $R$,
  $t_1$, \dots, $t_n$ is uniquely determined.
\item $A \ident \eq[t_1][t_2]$ where $t_1$ and $t_2$ are uniquely
  determined terms.
\end{enumerate}
\end{prop}

\begin{proof}
Exercise.
\end{proof}

\begin{prob}
Prove \olref[fol][syn][unq]{prop:unique-atomic} (Hint: Formulate and
prove a version of \olref[fol][syn][unq]{lem:no-prefix} for terms.)
\end{prob}

\begin{prop}[Unique Readability]
Every \olup{formula} satisfies one, and only one of the following conditions.
\begin{enumerate}
\item $A$ is atomic.

\tagitem{prvNot}{$A$ is of the form $\lnot B$.}{}

\tagitem{prvAnd}{$A$ is of the form $(B \land C)$.}{}

\tagitem{prvOr}{$A$ is of the form $(B \lor C)$.}{}

\tagitem{prvIf}{$A$ is of the form $(B \lif C)$.}{}

\tagitem{prvIff}{$A$ is of the form $(B \liff C)$.}{}

\tagitem{prvAll}{$A$ is of the form $\lforall[x][B]$.}{}

\tagitem{prvEx}{$A$ is of the form $\lexists[x][B]$.}{}
\end{enumerate}
Moreover, in each case $B$, or $B$ and $C$, are uniquely
determined.  This means that, e.g., there are no different pairs $B$,
$C$ and $B'$, $C'$ so that $A$ is both of the form 
\iftag{prvIf}{$(B \lif C)$ and $(B' \lif C')$.}{
    \iftag{prvOr}{$(B \lor C)$ and $(B' \lor C')$.}{
      \iftag{prvAnd}{$(B \land C)$ and $(B' \land C')$.}{}}}
\end{prop}

\begin{proof}
The formation rules require that if \olup{formula} is not atomic, it
must start with an opening parenthesis~(, \iftag{prvNot}{$\lnot$,}{}
or a quantifier. On the other hand, every \olup{formula} that starts with
one of the following symbols must be atomic: \olup{predicate}, \olup{function}, \olup{constant}\iftag{prvFalse}{, $\lfalse$}{}\iftag{prvTrue}{, $\ltrue$}{}.

So we really only have to show that if $A$ is of the form $(B \ast
C)$ and also of the form $(B' \mathbin{\ast'} C')$, then $B \ident
B'$, $C \ident C'$, and $\ast = {\ast'}$.

So suppose both $A \ident (B \ast C)$ and $A \ident (B'
\mathbin{\ast'} C')$.  Then either $B \ident B'$ or not.  If it is,
clearly $\ast = {\ast'}$ and $C \ident C'$, since they then are
substrings of $A$ that begin in the same place and are of the same
length.  The other case is $B \not\ident B'$.  Since $B$ and
$B'$ are both substrings of $A$ that begin at the same place, one
must be a proper prefix of the other.  But this is impossible by
\olref{lem:no-prefix}.
\end{proof}



\end{document}
